\documentclass[a4paper,11pt]{article}
\usepackage[utf8]{inputenc}
\usepackage[T1]{fontenc}
\usepackage{lmodern,amsmath,amssymb,graphicx,booktabs,longtable,array}
\usepackage[margin=24mm]{geometry}
\usepackage{microtype}
\usepackage[colorlinks=true,linkcolor=blue,citecolor=blue,urlcolor=blue]{hyperref}
\setlength{\emergencystretch}{2em}
\newcommand{\dd}{\mathrm{d}}
\newcommand{\ii}{\mathrm{i}}
\title{Radiative cancellations and localized linear modes\\in three-dimensional sextic Q-balls}
\author{Finn Malte Hinrichsen\\\small Hamburg, Germany}
\date{4 October 2026 --- Draft 0.45}
\begin{document}
\maketitle
\begin{abstract}
We study localized linear excitations of a complex scalar Q-ball in three
spatial dimensions with potential $\mathcal U(S)=S-S^2+S^3/2$.
The linearized problem couples two frequency sidebands. In a region with
one propagating and one evanescent exterior channel, radiation can cancel
at particular parameter values. We assemble computer-assisted existence
arguments for two radial modes and one dipole mode, together with a
numerically observed sequence of fifteen radial candidate locations.
The existence statements concern the infinite radial domain, using
analytic exterior estimates and ball-arithmetic interior enclosures;
they are not inferred from small finite-box linewidths. A thin-wall
phase-matching description, partly calibrated, describes the observed near-uniform
spacing in inverse detuning. Sequential and reportedly author-blinded comparisons
test predictions for further linewidth minima, but do not certify the complete sequence
or an infinite ladder. We distinguish these linear results from nonlinear
confinement: a direct expansion produces second-harmonic sources in open
channels, including monopole and quadrupole channels for a real dipole
excitation. A numerical forced-response calculation supports nonzero
outgoing amplitudes for an approximate dipole input. A complete monopole
error budget additionally separates one second-harmonic coefficient from
zero, conditional on the stated exact-mode and analytic enclosure
hypotheses. This supplies an obstruction to a specified uniformly
localized smooth continuation of a real dipole; circular polarization
remains outside this monopole conclusion. The
manuscript makes no claim of linear or nonlinear stability, a particle model, or an
experimental realization.
\end{abstract}
\noindent\textbf{Status.} This is a complete working draft, not a submitted or
journal-refereed article. The proof components have project-internal
cross-house reviews; the assembled manuscript has its own narrower
editorial checks. A second independent interval implementation is not
provided. Numerical values are dimensionless throughout.
\tableofcontents

\section{Question and scope}\label{sec:scope}
A long-lived oscillation and an exactly localized linear mode are different
objects. A numerically small resonance width may result from a cancellation,
a reflecting boundary, or insufficient resolution. Conversely, even a
localized linear mode need not extend to a localized finite-amplitude
oscillation: nonlinear harmonics can couple to the continuum.
Our objective is to separate these statements within a single specified
scalar field model. We report three concrete existence enclosures,
a broader numerical sequence, and an approximate mechanism for its spacing.
The sequence index $n$ is a continuation label, not a separately certified
node count or a statement that no intermediate modes exist.

A limited mechanical analogy is a bell with several vibration patterns:
exciting a particular pattern differs from exciting an arbitrary disturbance.
Here the additional ingredient is cancellation of the outgoing amplitude
in an available radiation channel, a general interference mechanism
discussed in Ref.~\cite{FriedrichWintgen1985}. The cancellation can be lost
under perturbations that do not preserve its condition. The analogy concerns
linear modes; it implies neither undamped finite-amplitude ringing nor
immunity to environmental disturbances.

The central contribution claimed here is this combination of model-specific
linear existence enclosures and numerically tested spacing structure.
We do not claim that embedded modes, interference suppression, Q-ball
excitations, or thin-wall approximations are new concepts. No priority claim
is made for the general mechanism. The nonlinear analysis is retained as
a limit on interpretation, not used to advertise stable particles.
\subsection{Relation to Q-ball excitations and continuum trapping}
\label{sec:related-work}

Q-balls belong to the established theory of localized scalar configurations
carrying a conserved global charge~\cite{Coleman1985}. Their perturbations
and nonlinear interactions have an extensive history. Battye and Sutcliffe
simulate collisions, charge transfer and fission in one, two and three
spatial dimensions in the same classical sextic model, after unit
conversion~\cite{BattyeSutcliffe2000}. Their kinetic normalization is
$\tfrac12|\partial\phi|^2$ and their potential is $2\mathcal U$;
$x=2x_{\mathrm{BS}}$, $t=2t_{\mathrm{BS}}$ give $\omega=\omega_{\mathrm{BS}}/2$.
Integrated energies and charges must also be converted in the chosen
spatial dimension. Their slowly decaying one-dimensional breathing
example is a dynamical precursor, not a certificate of an embedded linear
mode. The model itself is not new here. Its perturbations
require a coupled frequency problem. Kovtun, Nugaev, and Shkerin formulate
the complex-field linearization with conjugate coupling and the two
frequencies $\omega\pm\rho$~\cite{Kovtun2018}. Their three-dimensional
partial-wave equations provide a direct structural comparison, while the
vibrational modes selected in their piecewise-parabolic example satisfy
$\omega+\rho<1$. That calculation does not establish the embedded modes
of the sextic potential studied here.

Azatov, Ho, and Khalil study three-dimensional scattering in the same
sextic model family~\cite{Azatov2024}. With their $\phi_Q=\sqrt2 f$ and
$g=\beta=1/2$, their linear coefficients $U,W$ equal our $D,C$.
Conjugating their carrier convention and reducing the radial functions
maps their Appendix~B partial-wave equations to
Eqs.~\eqref{eq:radialA}--\eqref{eq:radialB}.
The radial operator is established prior work; the question here is the
existence of the particular localized solutions enclosed below.

The same work also treats Friedberg--Lee--Sirlin Q-balls containing a
complex and a real scalar. Its Section~3 derives the coupled charged
sidebands and neutral channel, their asymptotic thresholds, and scattering
relations~\cite{Azatov2024}. Thus an additional real field and a
three-channel linearization are not new ingredients. These scattering
calculations do not provide the particular infinite-domain embedded-mode
enclosures examined here.

Zhang, Zhou, and Zhu obtain analytic scattering solutions for stepwise
approximations to Q-balls of the same sextic
family~\cite{ZhangZhouZhu2025}. With their $g=\beta$ and
$f_0^2=1/(2\beta)$, their interior wavenumber $\sqrt{-\rho_1}$,
Eq.~(101), equals $k_{\mathrm{in}}$ in Eq.~\eqref{eq:interior}.
Their propagating scattering regime requires both exterior channels to
be open. The common interior dispersion and thin-wall phase structure
are therefore established ingredients; they do not themselves establish
cancellation of the remaining open-channel amplitude in the
one-open-channel regime considered here.

Ciurla, Dorey, Roma\'nczukiewicz, and Shnir study a sextic Q-ball model in
$1+1$ dimensions~\cite{Ciurla2024}. Their separation of bound,
half-propagating, and quasinormal modes is particularly relevant: in the
interval $1-\omega<\rho<1+\omega$, one asymptotic channel propagates and
the other is evanescent. They exhibit a small but nonzero radiation tail
associated with a long-lived quasinormal mode and investigate nonlinear
relaxation. This distinguishes a narrow resonance from the exact
vanishing of the propagating tail required here. Their one-dimensional
examples do not supply the spherical profile, origin conditions, or
infinite-domain enclosure used in the present three-dimensional problem.

Evslin et al. analyze paired Floquet modes of small-amplitude Q-balls in
$1+1$ dimensions~\cite{Evslin2026}. Their counterrotating analysis relates
a discrete component coupled to a propagating partner to a Feshbach-type
quasinormal mode. The small-amplitude approximation and one-dimensional
profile do not supply the full spherical sextic problem or the vanishing
open-channel amplitude examined here.

Continuum trapping through channel interference is also an established
mechanism. Friedrich and Wintgen relate interfering resonances to a
vanishing decay width~\cite{FriedrichWintgen1985}. A rigorous related
construction by Yu and Lu treats three Schr\"odinger channels on a
half-line, with one open channel, a crossing of two uncoupled bound
levels, and suitable nonzero radiation couplings~\cite{YuLu2025}.
Their small-coupling theorem concerns a potential that is exactly constant
outside a finite interval. It is not directly applicable to the
frequency-dependent two-channel Q-ball pencil, whose coefficients derive
from a nonlinear profile and retain an infinite tail. We therefore use
these works as comparisons, without assigning the specific
Friedrich--Wintgen mechanism to our modes solely from a vanishing
radiation amplitude.

Resonant total reflection through coupled propagating and evanescent
channels is established for discrete breathers~\cite{Flach2003},
optical solitons~\cite{Flach2005}, and spinor-condensate domain
walls~\cite{Watabe2012}. These are scattering solutions with an incident
wave; localization of a closed-channel component is not by itself an
embedded eigenmode of the full problem. Nor does the existence of open
and closed channels suffice for reflection: the fully coherent homogeneous
optical example remains transparent~\cite{Flach2005}.
These comparisons motivate a broader channel-interference interpretation,
without making an antiparticle branch a universal prerequisite or
establishing a necessary condition on Krein signatures.

Accumulation of embedded solitons also predates this work. Malomed et al.
derive a WKB/Bohr--Sommerfeld description in a quadratic optical model,
with a predicted scaling $\varepsilon_n\sim n^{-6/5}$~\cite{Malomed2005}.
Their matching is explicitly nonrigorous and its leading prefactor differs
from their numerical fit. Their objects are nonlinear embedded solitons;
ours are linear internal modes on changing Q-ball backgrounds. Neither
the scaling parameters nor their exponents can be identified merely from
the existence of a sequence of radiation cancellations.

Soffer and Weinstein prove nonlinear radiation damping under a positive
Fermi-golden-rule hypothesis for a real cubic Klein--Gordon
model~\cite{SofferWeinstein1999}. Its initial discrete mode lies below the
continuum and its third harmonic radiates. Their theorem does not directly
cover our embedded mode, rotating background, or second-harmonic source.
Our conditional continuation obstruction is not a dynamical decay theorem.

Suppressed radiation also occurs for real-field oscillons. Zhang et al.
report multiple decay-rate minima for the $\phi^6$ potential in
Section~7.5~\cite{ZhangOscillons2020}. Their analysis concerns long-lived
nonlinear configurations with radiative decay, rather than an enclosure
of an exact embedded eigenmode. The authors also identify limitations in
lifetime estimates near decay-rate minima. A small radiation signal or
a long-lived oscillation is therefore not, by itself, the criterion
used for the linear localization claims here.

An especially close comparison for the nonlinear radiation problem is
provided by Inagaki and Murakami~\cite{InagakiMurakami2026}.
They numerically resolve five zeros of the second-harmonic source overlap
for an internal mode of a thermal critical bubble, including three on
the thinner-wall side. The fundamental mode is below the continuum;
it is the leading second-harmonic radiation amplitude that vanishes.
At a selected zero, third-harmonic radiation remains and controls the
leading small-amplitude decay. These results neither constitute a linear
embedded eigenmode nor establish a radiation-free nonlinear solution.
They also do not establish an infinite sequence or its asymptotic spacing.
Our linear localization claim concerns the fundamental mode itself;
this distinction does not by itself establish priority.

The contribution examined here is the existence of particular localized
solutions of that full linearized Q-ball problem at the stated parameter
boxes, with regular origin conditions and controlled infinite tails.
The literature above does not replace those checks, and this comparison
does not establish priority. A finite collection of certified points also
does not imply an infinite ladder, or that modes on distinct backgrounds
belong to one fixed Q-ball spectrum. Finally, linear localization alone
does not establish an exactly localized nonlinear oscillation of finite
amplitude or nonlinear stability.


\section{Field model and two-channel linearization}\label{sec:model}
With metric signature $(+,-,-,-)$, define
\begin{equation}\label{eq:action}
\mathcal L=\partial_\mu\phi^*\partial^\mu\phi-\mathcal U(|\phi|^2),
\qquad \mathcal U(S)=S-S^2+\beta S^3,\qquad \beta=\tfrac12.
\end{equation}
The free scalar mass is one. A spherical stationary background is
$\phi_0=e^{\ii\omega t}f(r)$ with real $f$, satisfying
\begin{equation}\label{eq:profile}
f''+\frac2r f'=(1-\omega^2)f-2f^3+3\beta f^5,
\quad f'(0)=0,\quad f(r)\longrightarrow0.
\end{equation}
For definiteness the charge sign is chosen so that
$Q=2\omega\int f^2\dd^3x>0$ when $\omega>0$.
The thin-wall frequency threshold is
$\omega_c^2=\min_{S>0}\mathcal U(S)/S=1-1/(4\beta)=1/2$.

Use a real orthonormal spherical-harmonic basis and write
\begin{equation}\label{eq:ansatz}
\delta\phi=e^{\ii\omega t}
 [a(r)e^{\ii\rho t}+b(r)e^{-\ii\rho t}]Y_{lm}(\hat x),
\qquad A=ra,\quad B=rb,
\end{equation}
for real-frequency, real-radial-function modes. In a complex harmonic basis
the second sideband instead carries $Y_{lm}^*$.
Let $S=f^2$, $D=1-4S+9\beta S^2$, and $C=-2S+6\beta S^2$.
The radial equations are
\begin{align}
 A''&=\left[D+\frac{l(l+1)}{r^2}-(\omega+\rho)^2\right]A+CB,
 \label{eq:radialA}\\
 B''&=CA+\left[D+\frac{l(l+1)}{r^2}-(\omega-\rho)^2\right]B.
 \label{eq:radialB}
\end{align}
Regular solutions satisfy $A,B=O(r^{l+1})$ at the origin.
In the range $1-\omega<\rho<1+\omega$, define
\begin{equation}\label{eq:channels}
k^2=(\omega+\rho)^2-1>0,\qquad
\kappa_c^2=1-(\omega-\rho)^2>0,\qquad
\kappa_0^2=1-\omega^2>0.
\end{equation}
The $A$ channel is open and the $B$ channel closed. A localized mode
requires cancellation of both oscillatory exterior components of $A$
and exclusion of the growing component of $B$.
Throughout, ``embedded'' means a localized solution in this explicitly
open-channel regime. A separate operator-domain and essential-spectrum
theorem for the full time-evolution generator is not asserted.

Spatial localization does not require compact support. A scalar component
of laboratory frequency $E$ with $|E|<1$ has the vacuum decay scale
$\kappa=\sqrt{1-E^2}$ and free decaying asymptotic proportional to
$e^{-\kappa r}/r$, with angular-momentum corrections. The threshold
$|E|=1$ requires separate treatment. For $|E|>1$ a propagating channel
exists, but its oscillatory exterior coefficients vanish for the localized
linear modes considered here; exponentially forced tails can remain in
the coupled problem. Infinite spatial support implies neither an
unsuppressed long-range force nor a universal interaction. A stationary
tail is also distinct from the causal propagation of a new disturbance.

\subsection{Spatial mode shapes and phase-resolved reconstruction}\label{sec:mode-figures}
The three-dimensional shape is obtained from the computed radial functions
and the specified spherical harmonic, not from a point-particle picture.
For real radial components and the positive-carrier convention of this paper,
we display the rotating-frame perturbation per unit small amplitude,
\begin{equation}
 \psi_1(\mathbf{x},t)=Y_{l0}(\widehat{\mathbf{x}})
 \left[\frac{A(r)}r e^{i\rho t}+\frac{B(r)}r e^{-i\rho t}\right].
 \label{eq:mode-render}
\end{equation}
The full first-order field is $e^{i\omega t}(f+\epsilon\psi_1)$.
Thus a stationary background density does not mean that the complex field
is static. A localized nonzero oscillation can have zero linear radiation;
this is distinct from removing the oscillation itself.

Figures~\ref{fig:radial-mode} and~\ref{fig:dipole-phases} use stored numerical
reconstructions near C0 and T2. They are not pointwise interval-certified
eigenfunctions and are not a nonlinear three-dimensional time evolution.
The T2 arrays were produced in the conjugate carrier convention; the time
factors in Eq.~\eqref{eq:mode-render} conjugate that convention consistently.
Interpolation is used only for display. All lengths are dimensionless model
coordinates; no particle diameter or laboratory length is assigned.

\begin{figure}[htbp]\centering
\includegraphics[width=.97\linewidth]{radial-mode.pdf}
\caption{C0 spherical mode at phase zero. Left: cutaway surfaces of constant
$|\psi_1|$ at25\%,50\% and75\% of its displayed maximum; color identifies
the amplitude level. Removing a quadrant exposes the radial shells
and is only a drawing choice. Right: the background and the radial mode,
separately normalized for comparison. The shells are level sets, not material
surfaces. These are numerical reconstructions associated with C0, not new
existence certificates.}\label{fig:radial-mode}
\end{figure}

\begin{figure}[htbp]\centering
\includegraphics[width=.97\linewidth]{dipole-phases.pdf}
\caption{T2 with a real $Y_{10}=\sqrt{3/(4\pi)}z/r$ angular factor.
Top: three phases of the linear mode at a common amplitude-isosurface level.
Bottom: the corresponding $y=0$ sections; phase is masked below1.5\% of the
common displayed peak, because phase is undefined at a node. Color encodes
$\arg\psi_1$ in the rotating frame, with a fixed common gauge and cyclic scale;
it is neither a vector direction nor an additional charge. The opposite
lobes are in antiphase. The changing quadratures follow the two computed
radial sidebands and their analytic time factors, rather than a fresh
nonlinear simulation.}\label{fig:dipole-phases}
\end{figure}

A spatial node of $\psi_1$ is not the BIC cancellation condition. The latter
is a condition on the frequency and background parameters. To visualize
this distinction, Fig.~\ref{fig:bic-winding} shows an archived C0 search using
a normalized backward-Jost origin residual $F=(F_1,F_2)$.
Its complex representation $F_1+iF_2$ is a device for tracking the winding
of two real boundary residuals, not the phase of the physical field.
The finite-radius numerical diagnostic predates and does not replace the
existence enclosure in Section~\ref{sec:proof-model}.

\begin{figure}[htbp]\centering
\includegraphics[width=.97\linewidth]{bic-winding.pdf}
\caption{The C0 cancellation in parameter space. Left: two sampled parameter
loops lifted by $|F|$, colored by $\arg(F_1+iF_2)$; the star marks the numerical
root. The reference coordinates are $\rho_{\mathrm{ref}}=1.744625038252254$ and
$\omega^2_{\mathrm{ref}}=0.7976953125$. Right: their residual images and an offset
control, with numerical windings$+1,+1,0$. Lines connect stored samples;
no two-dimensional residual surface has been inferred from these loops.
The historical calculation uses outer radius$44$. The plotted residual is
neither radiated power nor an absolute outgoing-wave amplitude, and the
third plotting coordinate is not a third physical-space dimension.}
\label{fig:bic-winding}
\end{figure}


\section{Evidence hierarchy and numerical sequence}\label{sec:sequence}
We separate three levels: a width minimum is a numerical candidate;
a resolved winding test of a specified two-component matching map gives
additional numerical evidence; an interval enclosure with analytic tails
supports an infinite-domain existence statement. These labels cannot be
interchanged. A complex pole search and a real matching zero are also
not identical numerical tasks.

% Table revised in v0.10 from independently audited primary records.
% Sources: DATENPAKET.json e34aa7f14712736b5c0a4d9bc69014c58e8fcc23859ad229188985f3da214a4f
% C0: ZERT-A2-L40.json 7fd55b3f494333fb86ce8ea4152a6c7908462df1fa3b403ae2d172e6053d0798
% C1: ZERT-T1-A-L32.json 97e13cb196318114d4ad1dba5311ad52a4881aaa14ef38ed66d98fe02d065177
% Full paths and provenance in data-audit/SOURCE-SHA256SUMS.txt and AUDIT-B2-B10-B11-B13.txt.
Table~\ref{tab:ladder} distinguishes the continuation locations by their
estimator and evidence class. Rows 1 and 2 use the subsequent existence
certificates C0 and C1. Rows 6--10 now use the subsequent R13 calculation;
the other numerical rows retain their archived data-package records.
Rows 3 and 4 have resolved \emph{numerical} winding
tests in the truncated radial problem, not interval-certified winding
enclosures. Row 5 has sign-change evidence, but its sampled winding test
did not meet the resolution criterion. Historically, for indices 7--9, pole-width
interpolation replaced a problematic direct sign calculation with strong
interior growth. For indices 7--9 and 11--15, locations were reconstructed
from short linewidth samples using signs selected to match neighboring
slopes; those archived records did not report resolved numerical BIC
winding tests. The old row-10 reference was an unrefined continuation
estimate, distinct from the observed width depression at the scan edge.

R13 adds real matching sign changes and resolved numerical windings for
$n=6,7,8,9,10$ in the truncated radial problem. On both radial steps
$h=0.04$ and $0.02$, the windings for $n=7,8,9,10$ are respectively
$-1,+1,-1,+1$, with maximum sampled phase increments below $0.301$
radians. The $n=1$ control agrees with the certified location to its
prescribed tolerance; the $n=6$ control gives approximately
$\omega^2=0.569360$ and a resolved positive winding. This is new
numerical evidence, not additional interval-certified existence results.
For $n\geq11$, the earlier linewidth evidence retains its previous status.

The R13 locations are linearly interpolated between $\omega^2$ samples
separated by $2.5\times10^{-4}$. Agreement between the two radial steps
within $2.3\times10^{-8}$ is therefore not a location error bound of that
size. The report's interpolation-uncertainty estimate of a few $10^{-6}$
is not a validated enclosure. Signed-width reconstructions differ from
these locations by at most $2.2\times10^{-6}$, but share profiles,
equations and implementation with the matching calculation; this is not
an independent replication. Candidate interior growth and window-maximum
growth are different diagnostics, and the adopted growth threshold is
a numerical convention, not an analytic error bound.
% R13 ERGEBNIS.md SHA256 2e428ccfbff1c26623e4b83caa3d8adee74b0f99467a69b2e6e9d6814ef0e037
\begin{table}[htbp]\centering\small
\caption{Radial sequence at $\beta=1/2$. Approximate candidate locations
are rounded from the indicated archived records or new R13 calculation;
displayed digits
do not specify confidence intervals or validated location errors.
For C0 and C1, the full enclosures are given in
Section~\ref{sec:certified-existence}.}
\label{tab:ladder}
\begin{tabular}{rll}\toprule
$n$&$\omega^2$&Evidence and source record\\\midrule
1&$0.797676787\ldots$&Existence enclosure C0\\
2&$0.685128904\ldots$&Existence enclosure C1\\
3&$\sim0.631450$&Resolved numerical winding; A03\\
4&$\sim0.601421$&Resolved numerical winding; A04\\
5&$\sim0.582417$&Sign change / curve refinement; A05\\
6&$\sim0.569360$&Resolved numerical winding; R13-K2\textsuperscript{a}\\
7&$\sim0.559848$&Resolved numerical winding; R13-n7\\
8&$\sim0.552619$&Resolved numerical winding; R13-n8\\
9&$\sim0.546940$&Resolved numerical winding; R13-n9\\
10&$\sim0.542364$&Resolved numerical winding; R13-n10\textsuperscript{b}\\
11&$\sim0.53860$&Linewidth reconstruction; A11\\
12&$\sim0.535449$&Linewidth reconstruction; A12\\
13&$\sim0.532772$&Linewidth reconstruction; A13\\
14&$\sim0.530470$&Linewidth reconstruction; A14\\
15&$\sim0.528469$&Linewidth reconstruction; A15\\\bottomrule
\end{tabular}
\par\smallskip
\begin{minipage}{.96\linewidth}\footnotesize
\textsuperscript{a}Historical A06 evidence, not the new R13 result:
two refinement runs ended near $0.569397$ and
$0.569379$, differing by about $1.8\times10^{-5}$. Neither reached the
prescribed $|\operatorname{Im}\rho|<10^{-8}$ refinement criterion;
convergence of the internal pole search does not establish convergence
of the location refinement. These two estimates do not define a validated
error interval.\par
\textsuperscript{b}Historical A10 evidence, not the new R13 result:
the scan showed a width depression at its edge $\omega^2=0.5425$;
the former reference $\sim0.5424$ was a continuation estimate,
not a previously refined pole-width minimum.
\end{minipage}
\end{table}

\begin{figure}[htbp]\centering
\includegraphics[width=.96\linewidth]{ladder.pdf}
\caption{Reported radial locations and inverse detuning
$1/(\omega^2-1/2)$. Filled symbols distinguish the two radial locations
with existence enclosures from the remaining numerical candidates.
No error bars are invented for rows lacking defensible intervals.
The plot uses the new R13 interpolated locations for $n=6$--$10$
($\sim0.569360$, $0.559848$, $0.552619$, $0.546940$, $0.542364$),
archived estimator values for the other numerical rows, and C0/C1 from
their certificates. R13 supplies numerical sign-change and winding
evidence, not new existence enclosures. Inverse detuning does not improve location
accuracy; this is neither a new fit nor a proof of accumulation.}
% Plot generated from the revised figure-data.json during the bound build.
\label{fig:ladder}
\end{figure}

\section{Sequential predictions and their limits}\label{sec:predictions}
Two prediction procedures were compared to the same target calculations.
The first extrapolated approximately constant inverse-detuning steps and
was updated after new results: first $n=11$, then $n=12,13$, then $n=14,15$.
It was not a single frozen five-point forecast. The second, a thin-wall
model fitted using indices 5,7,8,9, was reportedly blinded to later results
at the author level. Some target results were already known to the
coordinating team. This is information separation, not a claim of global
chronological preregistration or an independent numerical replication.
The frozen thin-wall forecast also included $n=10$ at
$0.542382\pm0.000020$. The new R13 interpolated location
$\sim0.542364$ is numerically compatible with this window. The old
$\sim0.5424$ reference was a continuation prediction, not a refined pole
minimum. The new comparison supplies neither a validated location error
nor a claim of exact zero linewidth. It also does not create retroactive
global blinding: the R13 code author's $n=7$ prediction followed a positive
smoke test, and is distinct from the coordinating team's earlier forecast.

\begin{table}[htbp]\centering\small
\caption{Reported locations and prediction windows. The $\pm$ values are
preassigned diagnostic windows, not statistical confidence intervals.
The two procedures used different acceptance rules.}\label{tab:forecast}
\begin{tabular}{rlll}\toprule
$n$&Numerical location&Sequential extrapolation&Thin-wall prediction\\\midrule
11&$\sim0.53860$&$0.53867\pm0.00030$&$0.538619\pm0.000025$\\
12&$0.535449$&$0.53544\pm0.00012$&$0.535468\pm0.000030$\\
13&$0.532772$&$0.53276\pm0.00012$&$0.532792\pm0.000035$\\
14&$0.530470$&$0.530460\pm0.000040$&$0.530490\pm0.000040$\\
15&$\sim0.528469$&$0.528458\pm0.000050$&$0.528489\pm0.000050$\\\bottomrule
\end{tabular}
\end{table}
The sequential location rule used one stated window; the thin-wall
comparison used its stated three-window rule. The handling of a failed
linewidth-threshold check also changed in the last sequential stage,
from failure to indeterminate, as documented before that stage.
All five reported comparisons passed their respective rules. The
thin-wall predictions lie systematically about $2\times10^{-5}$ above
the reconstructed locations. This is a visible model discrepancy, not
five independent random errors. The first sequential prediction lies
above its target, whereas the later ones lie slightly below.
Neither successful comparison establishes a zero linewidth.

\section{A thin-wall phase-matching description}\label{sec:thinwall}
Put $\epsilon=\omega^2-\omega_c^2$ and $S_c=1/(2\beta)$.
The leading thin-wall profile and radius in three dimensions are
\begin{equation}\label{eq:thinwall}
R_{\mathrm{tw}}\simeq\frac{1}{2\sqrt\beta\,\epsilon},\qquad
f^2(r)\simeq\frac{S_c}{1+\exp[(r-R_{\mathrm{tw}})/\sqrt\beta]}.
\end{equation}
These are asymptotic approximations, not substitutes for the certified
profiles. The leading radius relation is established prior work: with
$\Delta=m_\phi^2-\omega_0^2$, Heeck et al. use
$\kappa^2=(\omega^2-\omega_0^2)/\Delta$ and
$R^*=\sqrt\Delta R$, and obtain $R^*\simeq1/\kappa^2$
in their Eq.~(40)~\cite{Heeck2021}.
Here $\Delta=1/(4\beta)$, which gives the radius in
Eq.~\eqref{eq:thinwall}. Their inflection-point radius and our
half-density radius need not have identical subleading coefficients.
Diagonalizing the constant interior channel matrix gives the
propagating interior wavenumber
\begin{equation}\label{eq:interior}
k_{\mathrm{in}}^2=\omega^2+\rho^2-D_c+
 \sqrt{4\omega^2\rho^2+C_c^2},\qquad
D_c=1+\frac1{4\beta},\quad C_c=\frac1{2\beta}.
\end{equation}
At the stated interior amplitude, this is the same dispersion branch as
Eq.~(101) of Zhang, Zhou, and Zhu~\cite{ZhangZhouZhu2025}; the interior
wavenumber is an established ingredient, not a new result here.
A regular interior standing wave reaches the wall with phase
$\Psi=k_{\mathrm{in}}R_{\mathrm{tw}}$. The bare closed-channel wall approximation
supplies a localized state, while matching it to the open interior
wave permits a sign-changing radiation amplitude. A useful leading
parametrization is
\begin{equation}\label{eq:phase}
k_{\mathrm{in}}(\omega_n,\rho_n)R_{\mathrm{tw}}(\omega_n)
 \simeq\pi[n+\theta(\epsilon_n)],\qquad
\rho_n\simeq\omega_n+c(\epsilon_n).
\end{equation}
The radius and interior dispersion follow from the model approximations;
the wall phase $\theta$ and dressed level shift contained in $c$ remain
calibrated. A separate coupled planar-wall calculation now supplies a
numerical transmission zero, but does not determine both finite-curvature
functions independently. Consequently, agreement of the calibrated
phase rule with the sequence is not a parameter-free prediction.

For $\beta=1/2$ and $\omega^2=1/2$, the exact planar profile is
$S(x)=[1+\exp(\sqrt2 x)]^{-1}$. Solving the two-channel scattering
problem, with one open and one evanescent channel on each side, gives
\begin{equation}\label{eq:planarzero}
\rho_z\simeq1.52414976213,\qquad
k_{\mathrm{in}}(\rho_z)\simeq1.92332459129.
\end{equation}
An independently implemented Numerov boundary-value calculation confirmed
this zero before reading the comparison calculation. It used a continuously
phase-normalized transmission amplitude, actual complex-transmission
residuals, flux and reciprocity controls, and separate mesh and boundary
refinements. A conservative empirical frequency uncertainty is $10^{-9}$;
this is not a certified enclosure. Two finite frequency scans found no
additional zero, but do not prove uniqueness throughout the open-channel
window. The planar scattering state has an incident and reflected wave;
it is not itself a square-integrable embedded mode.

If $k_{\mathrm{in}}\to k_\infty>0$ and the wall phase has a finite limit,
Equation~\eqref{eq:phase} predicts, within the approximation,
\begin{equation}\label{eq:spacing}
\epsilon_n\sim\frac1{b_\infty n},\qquad
b_\infty=\frac{2\sqrt\beta\,\pi}{k_\infty}.
\end{equation}
From the rounded locations in Table~\ref{tab:ladder}, the inverse-detuning
steps for indices 12--15 are approximately 2.3043, 2.3053 and 2.3068;
these descriptive values do not specify accuracy beyond the input locations.
The calibrated model gives 2.298, while its bare
wall-state estimate gives 2.334. Substituting the independently computed
planar zero into Equation~\eqref{eq:spacing} instead gives
$b_\infty^{\mathrm{wall}}\simeq2.31000162857$, with a conservative empirical
numerical uncertainty of $2\times10^{-9}$. This wall-based candidate was not
fitted to ladder locations. Identifying it with the actual radial limiting
step still requires a controlled curvature and phase-matching limit.
This supports a useful spacing description;
it does not prove existence of infinitely many localized modes or identify
an exact limiting constant. The phase model also motivates alternating
signs of a matching Jacobian and quadratic linewidth minima, but neither
property independently replaces a tail-controlled existence proof.
The inverse-first-power scaling follows from the calibrated model, not
from an independently established asymptotic exponent. Exploratory free
power fits yielded exponents between $-0.76$ and $-0.88$; the reciprocal
fit improved the comparison metric only by factors of approximately
$1.2$--$2$, below the originally sought factor of ten. Finite-data agreement
therefore does not establish an exact asymptotic law.

A local numerical check also supports a planar transmission cancellation
for the modified potential $U(S)=S-S^2+S^3$ at
$\omega_{\min}=\sqrt{3}/2$. An independent Numerov implementation,
with the candidate frequency already known, gives
$\rho_z\simeq1.77345307180$ and $k_{\mathrm{in}}\simeq2.3994321234$.
The last grid change, $h=0.04$ to $0.02$, shifts $\rho_z$ by
$1.64\times10^{-10}$, while increasing the half-domain from $32$ to $40$
at the finest spacing changes it by $4.2\times10^{-14}$; these are
empirical sensitivities, not certified bounds. The corresponding thin-wall
spacing candidate is $2\pi/k_{\mathrm{in}}\simeq2.618613482$. This known-target
local replication neither counts all planar zeros nor establishes the
radial ladder limit for the modified potential.

An exploratory two-dimensional follow-up found further numerical
candidates near $\omega^2=0.529265$ and $0.525780$ after refining the
frequency sampling. The original decision rule remained inconclusive;
the resolved numerical winding for the former candidate and the
sign-change/linewidth evidence for the latter came from adaptively
specified follow-up runs. The last estimated inverse-detuning steps are
about $4.62$, compared with about $2.305$ in the three-dimensional sequence.
This is not a dimensional scaling theorem. A similar sampling limitation
in the earlier three-dimensional search remains a hypothesis.

\section{Computer-assisted existence results}
\label{sec:certified-existence}

We distinguish the three computer-assisted existence results below from
the numerically observed ladder. The proof establishes a localized mode
at a parameter pair inside each specified box. It does not certify every
observed radiation minimum, an infinite ladder, or a nonlinear
nonradiating solution.

\subsection{Model and precise statement}
\label{sec:proof-model}
Consider the complex Klein--Gordon equation in three spatial dimensions,
\begin{equation}
 \partial_t^2\phi-\Delta\phi+U'(|\phi|^2)\phi=0,
 \qquad U(S)=S-S^2+\tfrac12 S^3.
 \label{eq:proof-kg}
\end{equation}
A stationary Q-ball has the form $\phi_0=e^{i\omega t}f(r)$, where
\begin{equation}
 f''+\frac{2}{r}f'=(1-\omega^2)f-2f^3+\tfrac32 f^5,
 \qquad f(0)=a,\quad f'(0)=0.
 \label{eq:proof-profile}
\end{equation}
We use a real orthonormal basis $Y_{\ell m}$ of spherical harmonics and
write a linear perturbation as
\begin{equation}
 \delta\phi=e^{i\omega t}
 \left[\frac{A(r)}{r}e^{i\rho t}
       +\frac{B(r)}{r}e^{-i\rho t}\right]Y_{\ell m}.
 \label{eq:proof-mode}
\end{equation}
In a complex harmonic basis the second sideband instead contains
$Y_{\ell m}^*$. With $S=f^2$, substitution in the linearization of
\eqref{eq:proof-kg} gives
\begin{equation}
 \begin{aligned}
 A''&=\left[1-4S+\tfrac92 S^2-(\omega+\rho)^2
                  +\frac{\ell(\ell+1)}{r^2}\right]A+CB,\\
 B''&=\left[1-4S+\tfrac92 S^2-(\omega-\rho)^2
                  +\frac{\ell(\ell+1)}{r^2}\right]B+CA,
 \qquad C=-2S+3S^2.
 \end{aligned}
 \label{eq:proof-radial-system}
\end{equation}
Regularity requires $A,B=O(r^{\ell+1})$. Define
\begin{equation}
 \kappa_0=\sqrt{1-\omega^2},\qquad
 k=\sqrt{(\omega+\rho)^2-1},\qquad
 \kappa_c=\sqrt{1-(\omega-\rho)^2}.
 \label{eq:proof-channel-numbers}
\end{equation}

\noindent\textbf{Computer-assisted existence statement.}
For each row of Table~\ref{tab:proof-cases}, the corresponding certificate
encloses a quadruple $(a_*,c_*,\rho_*,\omega_*)$ with $c_*>0$ such that
\eqref{eq:proof-profile} has a positive global solution satisfying
$r e^{\kappa_0r}f(r)\to c_*$, and
\eqref{eq:proof-radial-system} has a nonzero real solution regular at the
origin with
\begin{equation}
 |A(r)|+|B(r)|\leq K_J e^{-\kappa_c r}\quad(r\geq L),
 \qquad e^{\kappa_c r}B(r)\longrightarrow1.
 \label{eq:proof-decay}
\end{equation}
The resulting spatial fields in \eqref{eq:proof-mode} are smooth and
exponentially localized. The certification is
conditional on the correctness of the stated interval-arithmetic
implementation and its underlying arithmetic libraries.

\begin{table}[htbp]
\centering
\small
\begin{tabular}{ccll}
\hline
Case & $\ell$ & Enclosure of $\omega_*^2$ & Enclosure of $\rho_*$ \\
\hline
C0 & 0 & $0.797676787111865191750\pm4.8\,10^{-22}$
       & $1.744617544837398735781\pm6.8\,10^{-22}$ \\
C1 & 0 & $0.6851289044582160933\pm2.5\,10^{-20}$
       & $1.6903565973281434568\pm4.7\,10^{-20}$ \\
C2 & 1 & $0.7544960137826372331\pm7.8\,10^{-20}$
       & $1.826342030181069956\pm3.0\,10^{-19}$ \\
\hline
\end{tabular}
\caption{The three certified parameter boxes, displayed as outward-rounded
decimal enclosures. The proof uses the full exact dyadic boxes in
$(a,c,\rho,\omega)$, not these two-coordinate summaries. C0 is BEWEIS-1;
C1 and C2 are respectively T1 and T2 of BEWEIS-2. Their numerical ladder
assignments are $(\ell,n)=(0,1),(0,2),(1,1)$; the labels $n$ are not
separately certified nodal counts.}
\label{tab:proof-cases}
\end{table}

For the enclosed parameter pair and its fixed radial angular channel, every
solution of \eqref{eq:proof-radial-system} that is square integrable on
$[L,\infty)$ is a scalar multiple of the constructed solution. This is
\emph{geometric radial simplicity}. It neither establishes algebraic
simplicity of a dynamical generator nor excludes Jordan chains. For C2,
the three-dimensional $\ell=1$ angular multiplet remains degenerate.
The translation mode at $\rho=0$ does not belong to the fixed nonzero
frequency considered here.

The term \emph{embedded} is used operationally: $k^2>0$ and
$\kappa_c^2>0$, so a free asymptotic channel is open, while the constructed
mode is exponentially localized and has no oscillatory radiation tail.
We do not identify the essential spectrum of a fully specified operator
pencil in this proof. The perturbation is periodic after removing
$e^{i\omega t}$; the two laboratory frequencies are $\omega\pm\rho$,
with no claim about their rational commensurability.

\subsection{Matching to an infinite exterior}
\label{sec:proof-tails}
The matching radius $L$ is not a truncation at which the potential is
discarded. Both the profile and the mode are continued to infinity by
Volterra constructions. For the profile set
$\widetilde u(r)=r e^{\kappa_0r}f(r)$ and
$q_f=-2f^2+\tfrac32 f^4$. Its tail satisfies
\begin{equation}
 \widetilde u(r)=c+
 \int_r^\infty\frac{1-e^{-2\kappa_0(s-r)}}{2\kappa_0}
                q_f(s)\widetilde u(s)\,ds.
 \label{eq:proof-profile-volterra}
\end{equation}
For a chosen bound $\overline f>0$, put
\begin{equation}
 \begin{aligned}
 K_P&=\frac{(2+\tfrac32\overline f^2)e^{-2\kappa_0L}}
                  {4\kappa_0^2L^2},\quad
 \eta=2K_Pc^3,\quad m=c+\eta,\\
 \Lambda&=\frac{(6+\tfrac{15}{2}\overline f^2)m^2e^{-2\kappa_0L}}
                     {4\kappa_0^2L^2}.
 \end{aligned}
 \label{eq:proof-profile-constants}
\end{equation}
The conditions $m e^{-\kappa_0L}/L\leq\overline f$,
$K_Pm^3\leq\eta$, $\Lambda<1$, and $c>\eta$ give a positive unique
fixed point in $\|\widetilde u-c\|_\infty\leq\eta$. Common interval
majorants verify these inequalities over the entire parameter box.
They also bound the corrections to both profile matching data.

For the mode, write $dV=-4f^2+\tfrac92 f^4$ and use
$\widetilde A=e^{\kappa_cr}A$, $\widetilde B=e^{\kappa_cr}B$.
In the $\ell=0$ case the normalized Volterra kernels are
\begin{equation}
 K_A(r,s)=e^{-\kappa_c(s-r)}\frac{\sin(k(s-r))}{k},\qquad
 K_B(r,s)=\frac{1-e^{-2\kappa_c(s-r)}}{2\kappa_c}.
 \label{eq:proof-l0-kernels}
\end{equation}
The inhomogeneous free term is $(0,1)$. For $\ell=1$ use
\begin{equation}
 \begin{aligned}
 j(x)&=\sin(x)/x-\cos(x),\quad y(x)=-\cos(x)/x-\sin(x),\\
 d(u)&=e^{-u}(1+1/u),\quad g(u)=e^u(1-1/u),\\
 G_A(r,s)&=\frac{j(kr)y(ks)-y(kr)j(ks)}{k},\\
 G_B(r,s)&=\frac{d(\kappa_cr)g(\kappa_cs)-g(\kappa_cr)d(\kappa_cs)}{2\kappa_c}.
 \end{aligned}
 \label{eq:proof-l1-kernels}
\end{equation}
Here $K_{A,B}=e^{-\kappa_c(s-r)}G_{A,B}$ and the free normalized
term is $(0,1+1/(\kappa_cr))$. For $kL,\kappa_cL\geq1$, valid bounds are
\begin{equation}
 \begin{aligned}
 |K_A|&\leq a_1=\frac{1+(kL)^{-2}}{k},&
 |K_B|&\leq b_1=\frac{b_0}{2\kappa_c},\quad b_0=1+(\kappa_cL)^{-1},\\
 |e^{-\kappa_c(s-r)}\partial_rG_A|&\leq a_2=\sqrt{1+(kL)^{-2}},&
 |(\partial_r-\kappa_c)K_B|&\leq b_2=1+(\kappa_cL)^{-1}
                                      +\tfrac12(\kappa_cL)^{-2}.
 \end{aligned}
 \label{eq:proof-l1-majorants}
\end{equation}
In particular, the last bound concerns the normalized \emph{physical}
derivative, not $\partial_rK_B$ alone. Since
$|dV|+|C|\leq6f^2+\tfrac{15}{2}f^4$ is integrable, ordered Volterra
iteration converges with factorial majorants. With
$Q\geq\int_L^\infty(|dV|+|C|)\,ds$ and $\mu=\max(a_1,b_1)$, the
$\ell=1$ correction bounds are $(a_1,b_1,a_2,b_2)\varepsilon$, where
$\varepsilon=b_0(e^{\mu Q}-1)/\mu$.
Uniform convergence also supplies the parameter continuity needed below.

\subsection{Finite-dimensional existence certificate}
\label{sec:proof-krawczyk}
The four unknowns are $z=(a,c,\rho,\omega)$. Rigorous forward integration
computes the profile and two regular radial solutions $R_1,R_2$ from
the origin. For $\ell=1$ the initial variables are
$A=r^2\alpha$, $B=r^2\beta$; their regular equations have derivative
term $4\alpha'/r$ or $4\beta'/r$. For a source $F$, the origin
integral equation has weight
$r^2\int_0^1\tau(1-\tau^3)F(\tau r)\,d\tau/3$, of mass $r^2/10$.
This avoids starting a singular $2/r^2$ equation with approximate data.
Subsequent interval Taylor steps use checked complex a priori enclosures
and Cauchy remainder bounds; parameter variations enclose the required
Jacobian over the whole box.

The matching map is $H=H_0+E$. Its first two rows match profile value
and derivative. The last two match the Wronski forms
$W(\Psi,R_i)=\Psi\mathbin{\cdot}R_i'-\Psi'\mathbin{\cdot}R_i$,
using free Jost data in $H_0$ and bounded tail corrections in $E$.
For $\ell=0$, these two Wronskis equal the origin values of $\Psi$.
For $\ell=1$, regular solutions behave as $r^2e_i$ and normalized
singular solutions as $r^{-1}e_i/3$; their Wronski pairing is
$\delta_{ij}$. Vanishing of both Wronskis therefore removes both
singular coefficients.

Let $Z=z_0+[-\delta,\delta]$, let $D$ enclose $DH_0(Z)$, and suppose
$|E_i|\leq\epsilon_i$ on $Z$. A fixed dyadic preconditioner $Y$ is used
to verify
\begin{equation}
 z_0-YH_0(z_0)+(I-YD)[-\delta,\delta]
             +|Y|[-\epsilon,\epsilon]\subset\operatorname{int}Z,
 \qquad
 \max_i\sum_j |I-YD|_{ij}\frac{\delta_j}{\delta_i}<1.
 \label{eq:proof-krawczyk-inclusion}
\end{equation}
The continuous map $z\mapsto z-YH(z)$ maps $Z$ into itself, so Brouwer's
theorem gives a fixed point. The weighted matrix bound makes $Y$
invertible, hence the fixed point satisfies $H(z_*)=0$. Differentiability
of $E$ is not required, and uniqueness of this parameter quadruple is
not asserted. Profile positivity follows from positive interior
enclosures, a minimum-principle argument on the remaining finite
interval, and the positive exterior fixed point.

Finally, $2\kappa_0>\kappa_c$, verified for all three boxes, permits
construction of two bounded oscillatory comparison solutions at infinity.
Their mutual Wronskian is nonzero. Together with the decaying Jost
solution they span a three-dimensional subspace of the four-dimensional
radial Cauchy-data space. Every square-integrable solution is
Wronski-orthogonal to this subspace, whose symplectic complement is
exactly the decaying line. This proves the geometric simplicity stated
above.

\subsection{Scope of verification}
\label{sec:proof-verification-limits}
The defining certificates use $L=40$ for C0, $L=32$ for C1, and $L=36$
for C2, with 256-bit ball arithmetic. Additional runs use larger matching
radii and 320-bit arithmetic. C2's originally specified repeat at $L=40$
failed the inclusion test. A subsequent run with an enlarged $c$-radius
passed after reevaluating the bounds on the new box; this is an adaptive
additional certificate, not an originally successful repeat. The C2
existence conclusion rests on its first successful certificate alone.

Project-internal cross-house mathematical and code reviews found no blocking error in
the reviewed chains; the stated scope and harmonic conventions include
their corrections. Review and higher-precision reruns do not constitute
a second independent implementation. Correctness of Python--FLINT,
Arb/Acb, and the recorded execution remains part of the computational
trust base. For C0 the simplicity inequality was recorded and separately
checked, but was not included in the historical aggregate success flag;
the newer C1/C2 drivers include it explicitly. Source and input versions
are tied to the archived certificates by hashes.

The validated arithmetic uses Arb/Acb midpoint--radius enclosures
\cite{Johansson2017Arb}, accessed through Python--FLINT
\cite{PythonFlintDocs,FlintSoftware}.
The finite-dimensional test is a Krawczyk-type preconditioned inclusion
\cite{Krawczyk1969}; for the separately bounded continuous tail correction,
existence follows from the explicit Brouwer argument in
Section~\ref{sec:proof-krawczyk}, without asserting uniqueness of the
parameter quadruple. These references identify the methodological
background; the integration, tail bounds and inclusion hypotheses remain
the specific checks of the archived certificates.

These results establish three linear localized modes in specified
parameter boxes. They do not establish nonlinear radiation cancellation,
linear or nonlinear stability, algebraic spectral simplicity, the absence of other
modes or parameter zeros, or the existence of every member of the
empirical ladder.

\section{Nonlinear sources and the limits of linear trapping}
\label{nl:section}

A localized eigenmode of the linearized equation need not remain
radiationless at finite amplitude. We distinguish an exact source
calculation, the kinematic opening of radiation channels, and a nonzero
outgoing response. Existing nonlinear evidence for radial location C0
is described below; it is not transferred to the additional modes T1
and T2 without evaluating their own sources.

\subsection{Quadratic expansion}
\label{nl:expansion}
We use the complex nonlinear Klein--Gordon equation
\begin{equation}
 \phi_{tt}-\Delta\phi+U'(|\phi|^2)\phi=0,
 \qquad U(S)=S-S^2+\tfrac12 S^3,
 \qquad \phi=e^{i\omega t}(f+z),
 \label{nl:model}
\end{equation}
where the real background satisfies
$-\Delta f-\omega^2f+U'(f^2)f=0$. Setting $s=f^2$, the linear operator is
\begin{equation}
 \mathcal Lz=z_{tt}+2i\omega z_t-\Delta z
 +[U'(s)+sU''(s)-\omega^2]z+sU''(s)\overline z.
 \label{nl:linear}
\end{equation}
The quadratic term in the equation's left-hand side is exactly
\begin{align}
 G_2[z]&=fU''(s)(z^2+2|z|^2)
       +\tfrac12 f^3U'''(s)(z+\overline z)^2 \nonumber\\
       &=\mathcal A z^2+2\mathcal A|z|^2+\mathcal C\overline z^{\,2},
 &\mathcal A&=f\left(-2+\tfrac92f^2\right),
 \qquad \mathcal C=\tfrac32f^3.
 \label{nl:quadratic}
\end{align}
Consequently $z=\epsilon z_1+\epsilon^2z_2+\cdots$ gives
$\mathcal Lz_1=0$ and $\mathcal Lz_2=-G_2[z_1]$. In particular, the
forcing has the opposite sign to the Taylor term in
Eq.~\eqref{nl:quadratic}.

For a real spherical harmonic $Y$ and real radial functions $a,b$, write
\begin{equation}
 z_1=Y\bigl(ae^{i\rho t}+be^{-i\rho t}\bigr).
\end{equation}
Then
\begin{align}
 G_2[z_1]&=Y^2\bigl(Q_+e^{2i\rho t}+Q_0+Q_-e^{-2i\rho t}\bigr),\\
 Q_+&=\mathcal A(a^2+2ab)+\mathcal Cb^2,\qquad
 Q_-=\mathcal A(b^2+2ab)+\mathcal Ca^2,\\
 Q_0&=2\mathcal A(a^2+b^2)+2(\mathcal A+\mathcal C)ab.
 \label{nl:coefficients}
\end{align}
These expressions introduce rotating-frame frequencies $0,\pm2\rho$,
or laboratory frequencies $\omega,\omega\pm2\rho$. The zero-frequency
term is static only in the rotating frame.

\subsection{Angular selection and forced channel problems}
\label{nl:selection}
For T1, the radial $l=0$ mode, $Y=Y_{00}$ and
$Y_{00}^2=Y_{00}/\sqrt{4\pi}$, so only $L=0$ occurs at this order.
For T2, choose a real normalized $l=1$ harmonic. After a spatial rotation
it equals $Y_{10}$, and
\begin{equation}
 Y_{10}^2=\frac{Y_{00}}{\sqrt{4\pi}}
          +\frac{Y_{20}}{\sqrt{5\pi}}.
 \label{nl:angular}
\end{equation}
Thus each of the three temporal source coefficients can drive $L=0$
and $L=2$; the antisymmetric $L=1$ product is absent. A different real
orientation rotates the quadrupole within the $L=2$ multiplet. This
standing real mode must not be confused with a rotating complex
$Y_{11}/\overline{Y}_{11}$ block: the latter has
$Y_{11}^2=\sqrt{3/(10\pi)}Y_{22}$ for its corresponding oscillatory product.

Let $c_L$ denote the coefficient in the relevant angular product and
$D_L=-\partial_r^2-2r^{-1}\partial_r+L(L+1)r^{-2}$. For the oscillatory
response $h_{+,L}e^{2i\rho t}+h_{-,L}e^{-2i\rho t}$ one must solve
the coupled system
\begin{equation}
 \begin{pmatrix}
 D_L+U'+sU''-(\omega+2\rho)^2&sU''\\
 sU''&D_L+U'+sU''-(\omega-2\rho)^2
 \end{pmatrix}
 \begin{pmatrix}h_{+,L}\\\overline{h}_{-,L}\end{pmatrix}
 =-c_L\begin{pmatrix}Q_+\\Q_-\end{pmatrix}.
 \label{nl:forced}
\end{equation}
Here $U',U''$ are evaluated at $s$; the conjugated second unknown matters
when imposing outgoing conditions. The real static correction obeys
$[D_L+U'+2sU''-\omega^2]h_{0,L}=-c_LQ_0$.
Gauge and charge constraints must also be fixed in constructing a
nonlinear continuation.

Both T1 and T2 have $|\omega|<1$ and $|\omega\pm2\rho|>1$.
For T2, the latter frequencies are approximately $4.52130$ and $-2.78407$.
The second harmonics are therefore open in the mass-one vacuum, whereas
the rotating-static correction is closed. The centrifugal term does
not change the asymptotic mass threshold. With the convention
$e^{iEt}$, an outgoing wave has radial phase
$e^{-i\operatorname{sgn}(E)\sqrt{E^2-1}\,r}$; a negative $E$ does not
mean negative outward energy transport.

An open channel and a nonzero local source do not establish radiation:
the outgoing source overlap can vanish. A localized continuation without
incoming radiation requires both open amplitudes to vanish in every
allowed $L$ block of Eq.~\eqref{nl:forced}. The T1 response remains open;
an approximate-input T2 calculation is reported below. A nonzero such amplitude would obstruct
a regular localized perturbative continuation with the specified leading
mode, but would not exclude all other nonlinear states. Frequency shifts
of order $\epsilon^2$ do not automatically cancel the second harmonic:
the carrier shift contributes a static background correction at this
order, while a shift of $\rho$ first acts on $\epsilon z_1$ at order
$\epsilon^3$.

\subsection{Complex dipole polarization and a conditional obstruction}
\label{nl:polarization}
The threefold dipole degeneracy allows more than rotations of a real
standing mode. Let $Y_i(\boldsymbol n)=\sqrt{3/(4\pi)}n_i$ be real
orthonormal Cartesian dipole harmonics and write
$A(\boldsymbol n)=\sum_i p_iY_i(\boldsymbol n)$, with
$\boldsymbol p\in\mathbb C^3$. For the same real radial eigenmode as above,
the allowed tangent has the linked sidebands
$z_1=a(r)A e^{i\rho t}+b(r)\overline A e^{-i\rho t}$.
This does not allow independently chosen polarizations of the two
sidebands or a mixture of distinct radial modes.

Set $\nu=\boldsymbol p^\dagger\boldsymbol p$ and
$\chi=\boldsymbol p^T\boldsymbol p$. The oscillatory quadratic forcing
is proportional to $A^2$, rather than $|A|^2$. Its orthogonal angular
projections satisfy the exact identities
\begin{equation}
 \|P_0(A^2)\|_{L^2(S^2)}^2=\frac{|\chi|^2}{4\pi},\qquad
 \|P_2(A^2)\|_{L^2(S^2)}^2
 =\frac{3\nu^2-|\chi|^2}{10\pi}\geq\frac{\nu^2}{5\pi}.
 \label{nl:polarization-norms}
\end{equation}
These follow by decomposing $\boldsymbol p\boldsymbol p^T$ into its trace
and symmetric traceless parts and integrating the fourth spherical
moment. Circular polarization, for example
$\boldsymbol p=(1,i,0)/\sqrt2$, has $\chi=0$ and hence no oscillatory
$L=0$ source, but its $L=2$ source does not vanish.

This source identity has a conditional implication for the outgoing
response. Define $C_{\mathrm{ref}}\in\mathbb C^2$ to be the \emph{exact}
radial $L=2,m=0$ outgoing coefficient for the normalized $Y_{10}$ tangent,
with the same radial normalization. Assuming the exact regular outgoing
problem exists uniquely, spherical symmetry makes its radial operator
independent of $m$. If $b_m=\langle Y_{2m},A^2\rangle$, linearity gives
\begin{equation}
 C_{2m}=\sqrt{5\pi}\,b_m C_{\mathrm{ref}},\qquad
 \sum_m\|C_{2m}\|^2
 =\frac{3\nu^2-|\chi|^2}{2}\|C_{\mathrm{ref}}\|^2
 \geq\nu^2\|C_{\mathrm{ref}}\|^2.
 \label{nl:polarization-response}
\end{equation}
The same identity holds separately for either radial component.
Conjugating the second computational channel to the physical sideband
permutes conjugate spherical harmonics and preserves these norm sums.
Thus an exact radial $L=2$ nonzero certificate would cover every nonzero
complex polarization in this one radial dipole eigenspace. The numerical
coefficients reported below do not provide that certificate.

For clarity, the corresponding nonlinear obstruction requires a precise
continuation class. Suppose a family is twice continuously differentiable
in amplitude as a $2\pi$-periodic rotating-frame profile in
$C^2(\mathbb T;H^4_\gamma(\mathbb R^3))$, where
$H^4_\gamma=\{w:e^{\gamma\sqrt{1+|x|^2}}w\in H^4\}$.
For the exact T2 background and tangent we may fix $\gamma=1/10$:
the inherited exterior decay rates are at least $0.495$ and $0.287$,
respectively. Their radial equations propagate this decay to derivatives
through order four; the regular origin expansions are even for the
background and of the form $r^2$ times even functions for the reduced
dipole components. Consequently the unreduced Cartesian fields are
smooth at the origin and belong to this same weighted space.
This membership statement does not assume that a nonlinear continuation
exists or certify a numerical Sobolev norm. Allow smooth
carrier and modulation frequency shifts and an additional static first
tangent, but no other first-order temporal harmonics. A nonzero exact
second-harmonic outgoing amplitude would contradict localization of the
second amplitude derivative. At this order the frequency-shift terms
are static or fundamental, so they cannot remove the second-harmonic
source. This excludes only that uniformly localized smooth continuation
class; it is not a dynamical instability or lifetime theorem.

\subsection{Numerical second-harmonic response of the dipole mode}
\label{nl:dipole-response}
We evaluated the T2 forcing in both $L=0$ and $L=2$ using separately
reconstructed numerical profiles near the certified point. The reduced
real mode components obey $\int_0^\infty(u^2+v^2)\,\dd r=1$, including
an approximate tail contribution, with $\int Y_{10}^2\,\dd\Omega=1$.
The solver uses the complex conjugate of our positive-carrier convention:
both computational unknowns have positive-$q$ outgoing Hankel conditions,
and the second unknown is conjugated to reconstruct the physical negative
sideband. The power is invariant under this convention change.

Three prescribed settings use exterior radii $48,64,64$, source spacings
$0.02,0.02,0.01$, and relative ODE tolerances
$2\times10^{-10},2\times10^{-10},2\times10^{-11}$. Two further solves at
the finest settings use the coarser reconstruction level from the same
input algorithm, with different exterior radius and tolerance. This is
an input-sensitivity comparison, not an independent implementation.
All four outgoing coefficients remain nonzero under these changes;
the largest observed relative complex-amplitude change is below
$7.8\times10^{-8}$. This sensitivity is not a rigorous error bound.

At the finest setting the harmonically averaged power coefficients are
\begin{center}
\begin{tabular}{ccc}
\toprule Angular sector&Positive sideband&Negative sideband\\\midrule
$L=0$&$2.4239\times10^{-7}$&$6.0359\times10^{-5}$\\
$L=2$&$4.4956\times10^{-7}$&$9.1587\times10^{-5}$\\\bottomrule
\end{tabular}
\end{center}
The leading power is $\epsilon^4$ times the sum of these coefficients in
the stated normalization. Both physical sidebands carry positive outward
energy with the stated outgoing conditions. These values neither determine a lifetime nor measure a
nonlinear time-evolution loss. They must not be compared directly with
coefficients using the earlier radial maximum-amplitude normalization.

This calculation supports second-harmonic leakage for the approximate T2
input, beyond merely finding open channels and nonzero local sources.
It does not certify leakage at the exact enclosed BIC. Float rounding
places some input parameters outside the tiny certified box; sampled
scaled Hermite-interpolant ODE defects have a sampled maximum of about
$5.2012\times10^{-6}$ (rounded upward to $5.3\times10^{-6}$).
These are midpoint samples on $0.01$-spaced grids extending to $r=80,96$,
excluding $|r-8|\le0.011$ around the splice, not uniform residual bounds.
Differences between reconstruction levels are much smaller. Omitted exterior
source and potential contributions were diagnosed, not converted into
propagated amplitude-error bounds in that calculation. The separate
validated monopole transfer in Section~\ref{nl:conditional-monopole}
addresses these errors using a fixed comparison function and the exact
mode's Jost normalization; it does not certify the numerical power table.

A subsequent fixed-coefficient reconstruction checked the source--current
identity for each angular sector. For the real symmetric forced operator,
$J=\operatorname{Im}(U^\dagger U')$ satisfies
$J'=\operatorname{Im}(F^\dagger U)$, where $U$ denotes the two reduced
response components and $F$ their forcing. Both computational channels
use positive-wavenumber outgoing waves. Keeping the saved matching
coefficients fixed, an independent volume quadrature agrees with
$\sum_j q_j|C_j|^2$ to relative residuals $9.6\times10^{-11}$ for $L=0$
and $1.3\times10^{-11}$ for $L=2$. This is a numerical consistency check
of the finite forced problem, not an energy-power measurement or a
validated error bound. It does not independently validate a shared source
formula. The change between the two volume quadratures is
$1.9\times10^{-9}$ or smaller relative to the respective outgoing current;
this observation is not a certified quadrature-error estimate.

% Conditional enclosure incorporated in v0.8; scope remains conditional.
% Numerical error budget does not remove the inherited analytical hypotheses.
\subsection{A conditional enclosure for the monopole response}
\label{nl:conditional-monopole}
For the real dipole tangent at T2, we assembled a complete error budget
for the $L=0$ second-harmonic forced problem. Here $L$ labels the
radiated angular sector, whereas the underlying localized linear mode
has $l=1$. The calculation uses the exact-mode Jost normalization,
with the evanescent reduced component satisfying
$\lim_{r\to\infty}e^{\kappa_c r}v(r)=1$.
It therefore has a different amplitude normalization from the numerical
power table in Section~\ref{nl:dipole-response}.

The error identity includes the inner and outer source differences,
operator differences, the residual of the fixed approximate response,
its actual matching jumps and origin contribution, and the exterior
boundary and coefficient-extraction errors. The exterior potential and
source are included through infinity. No measured small linewidth or
comparison between two numerical resolutions substitutes for these terms.

The following inherited hypotheses specify the exact problem to which
the numerical error budget applies; they are not new conclusions of
the scalar summation:
\begin{enumerate}
\item The T2 existence enclosure of Section~\ref{sec:certified-existence}
supplies one common exact parameter tuple
$z_*=(a_*,c_*,\rho_*,\omega_*)$ in its certified dyadic box, a positive
regular spherical profile $f_*$, and a nonzero real regular dipole mode.
All profile, mode and parameter-error bounds refer to this same tuple.
Neither uniqueness of the tuple nor the claim that every point of the
box is an exact root is assumed.

\item At that exact tuple, the decaying fundamental Jost solution belongs
to the span of the two regular origin solutions. This is the inherited
regular/Jost identification from the T2 Wronskian conditions, not an
inference from overlap of numerical residual intervals. Its reduced
components are $O(r^2)$ at the origin and obey the same normalization
$e^{\kappa_c r}v(r)\to1$, with
$\kappa_c=\sqrt{1-(\omega_*-\rho_*)^2}$.
The background and mode satisfy the analytic infinite-tail bounds:
for $r\ge36$ one may use
$|f_*(r)|\le(97/8)e^{-0.495r}/r$ and
$|u(r)|+|v(r)|\le1.097e^{-0.287r}$.

\item The exact quadratic $L=0$ source is the complete product source
of the real standing dipole in Eqs.~\eqref{nl:quadratic}--\eqref{nl:forced}.
The fixed nominal mode, source, response and comparison coefficient
use the same amplitude convention: the nominal unit-$L^2$ mode is divided
by the exact rational
$s_0=2558188928468177/2251799813685248$, and quadratic quantities by $s_0^2$.
This does not identify $s_0$ with the norm of the exact mode.
The computational carrier conjugation is retained throughout.

\item For the second harmonic, the exact reduced operator is
$-\partial_r^2-Q_*^2+V_*(r)$ with real symmetric, locally bounded
$V_*$ and
$Q_*=\operatorname{diag}(q_+,q_-)$,
$q_\pm=\sqrt{(\omega_*\pm2\rho_*)^2-1}$.
Both channels satisfy $5/2<q_\pm<23/5$.
The exterior Volterra bounds hold at these exact frequencies,
with the analytic potential and source tails retained through infinity.
At $R=64$, the exterior potential satisfies
$\int_R^\infty\|V_*(r)\|_2\,dr/(5/2)<1/2$ and the source is integrable.
Thus the exact outgoing reduction includes both its homogeneous maps
and its inhomogeneous source corrections. Together with regular
Dirichlet data at the origin, the positive outgoing matrix current
gives the unique forced solution and adjoint extraction problem.
This argument concerns two open second-harmonic channels, not the
fundamental embedded mode.

\item The analytic and validated input bounds used in the error identity
enclose the entire core and exterior, including exact-to-nominal
frequency and map changes. In Euclidean vector and induced matrix norms,
the inherited bounds
$\int_0^\infty\|V_*(r)\|_2\,dr<25/4$ and the Volterra estimates give
the adjoint amplitude weights $\|Z\|_2<26/5$ and $\|Z'\|_2<14$.
Source and operator differences, response residuals, origin and matching
jumps, exterior source terms, and coefficient-extraction rounding are
all bounded in this same convention, without dropping a tail or counting
one twice. The directed-arithmetic implementation and its interval
libraries remain part of the computational trust assumptions.
\end{enumerate}
Under these inherited hypotheses, the computed common coefficient-error
bound is

\begin{equation}
 \|C_* - C_{\mathrm{num}}\|_2
 \leq \mathcal E < 0.001023269269845280.
\end{equation}
The frozen comparison coefficient in the second computational channel
has the lower bound
\begin{equation}
 |C_{\mathrm{num},2}| > 0.001582591541472706.
\end{equation}
The directed arithmetic establishes, before rounding for display,
\begin{equation}
 |C_{*,2}| > 0.000559322271627427.
\end{equation}
This channel has laboratory frequency $\omega_*-2\rho_*<0$.
The solver uses a negative carrier and conjugates its second computational
component to recover that physical sideband. Conjugating the entire field
to the positive-carrier convention used here gives
$C_{-,\mathrm{paper}}=C_{*,2}$; all these operations preserve the modulus.
The first channel remains undecided by this common error bound; that
failure does not imply its coefficient vanishes.

This result supplies the nonzero-coefficient premise of the conditional
smooth-continuation argument for a real standing dipole. Its extension
to a complex dipole polarization is proportional to
$\boldsymbol p^T\boldsymbol p$ and therefore does not cover circular
polarization. It gives neither a nonlinear lifetime nor a dynamical
instability theorem. In particular, the bound is not a certified value
for any power coefficient in the preceding approximate-input table.


\subsection{Existing evidence for radial location C0}
\label{nl:evidence}
Separate computations already address the earlier $n=1$, $l=0$ mode
near $\omega^2=0.7976767871$ and $\rho=1.7446175448$. They concern the
same sextic field equation, not T1 or T2. A finite-amplitude radial
evolution compared the nonlinear field with its background and linear
response on two spatial grids. For amplitudes $0.01,0.02,0.04$, the
integrated quadratic flux of the nonlinear remainder increased by
factors approximately $15.97$ and $15.89$ upon doubling amplitude.
This supports fourth-order scaling over the tested range; that remainder
flux is not an exact measurement of the total lost core energy.

A subsequent harmonic diagnostic identified the dominant laboratory
sideband $E=\omega-2\rho\simeq-2.59611$ as outgoing at two detector
radii, on both grids and in both prescribed time windows. The weaker
$E=\omega+2\rho$ component was resolved but did not satisfy every
outgoing-wave criterion. Large residuals in the full field's harmonic
fit remained, so these observations do not establish a periodic
multifrequency solution.

A separately solved frequency-domain problem, using the same background
and mode data but no fitted amplitude from the time evolution, predicted
the dominant outgoing coefficient. At $\epsilon=0.02$, its predicted
amplitude was $5.1563\times10^{-5}$, compared with measured values
approximately $(5.03$--$5.10)\times10^{-5}$. The prescribed comparison
criteria passed. Shared input profiles make this a separate mechanism
test, not a wholly independent reconstruction of the background.

An interval calculation further separated one necessary overlap from
zero for an explicitly frozen finite operator on $[0,44]$:
\begin{equation}
 I_2\in[-0.16383470526432667,-0.10547329373338812].
 \label{nl:frozen}
\end{equation}
The corresponding $I_1$ interval contained zero, and the original test
requiring both exclusions remained failed. Nevertheless, one certified
nonzero overlap suffices to violate the necessary simultaneous condition
$I_1=I_2=0$ for that frozen problem. Its binary64 spline inputs are treated
as exact data, with a free exterior and zero source beyond radius $44$.
This is not a certificate that the exact continuum Q-ball and its BIC
have that overlap: a validated link to the continuum inputs and exterior
tail remains required. The linear existence certificates do not, by
themselves, supply that quantitative link.

If a leading second-order outgoing amplitude is nonzero, its power is
of order $\epsilon^4$. A lifetime law additionally requires a controlled
relation to mode energy and slow amplitude evolution, and separation of
initial adjustment from persistent radiation. Neither the finite-time
observations nor the frozen-operator certificate provides a general
nonlinear stability or lifetime theorem. Earlier real-field shell-model
normal forms describe a different model and supply no substitute for
these Q-ball source and overlap conditions.


\section{Interpretation, limitations and next discriminating tests}
The localized modes are excitations of an extended classical field
background. Nothing in this calculation identifies them with quarks,
spin-$1/2$ particles, gauge bosons, or dark matter. Angular degeneracy is
not a derivation of a local internal gauge symmetry. Likewise, geometric
simplicity of a radial solution space does not imply algebraic simplicity
of a frequency eigenvalue or nonlinear orbital stability.
The positive conserved form on a selected mode is neither a linear nor a
nonlinear stability theorem for the background or its full perturbation space.

\subsection{Environmental sensitivity and possible formation mechanisms}
Equation~\eqref{eq:action} contains a scalar field with global U(1)
symmetry, but no electromagnetic field or coupling to other particle
species. Its nonzero tail does not by itself imply interaction with
cosmic-microwave-background photons. Such a calculation would require
an explicit coupling and a dimensional scale.

On a fixed background, the conditional mode-projection identity separates
homogeneous scattering from prescribed local forcing: a homogeneous
solution with initially zero projection cannot populate it when the Green
boundary form vanishes, whereas a local source with nonzero overlap can.
The assumptions, forcing comparisons and source-work diagnostics are
retained in Appendix~\ref{app:driving}. They are distinct from the linear
existence enclosures and the nonlinear radiation obstruction above.

\paragraph{An additional radiation channel.}
A separate leading-order probe adds a neutral scalar through
$\mathcal L_{\mathrm{int}}=-g\chi|\phi|^2$. On three fixed radial profiles,
the finite-grid form factor
$I(k)=\int_0^\infty f(a+b)j_0(kr)r^2\,dr$,
$k=\sqrt{\rho^2-m_\chi^2}$, had respectively one, one and two detected
sign-changing zeros in the sampled interval $0<k<\rho$.
These are candidate Born-amplitude cancellations, not a complete count
of zeros of the exact integral. The coupling already sources a static
neutral field through $(-\Delta+m_\chi^2)\chi_0=-gf^2$; thus the
unchanged Q-ball with $\chi=0$ is not a solution of the coupled theory.
Background and mode backreaction must be included before claiming a
coupled bound state. Closing this channel with a strict threshold gap
removes propagation at this harmonic, but does not by itself preserve
the old channel cancellation or prohibit radiation at higher harmonics.

The exploratory concentration tests use a different, nonstationary initial
background. Appendix~\ref{app:formation} retains the non-dilution bound,
angular response, exterior diagnostic and constrained transport comparisons,
including their preparation dependence and finite-time limits. None
establishes formation of a stationary Q-ball or binding. These tests do
not supply additional evidence for the certified BIC backgrounds.
The appendix also records a later spherically symmetric relaxation study:
two individual Gaussian starts passed its charge--frequency diagnostic,
but its predeclared combined formation criterion was not met. That
limited observation does not resolve the broader formation question.

\subsection{Exploratory sensitivity to the numerical lattice}

Exploratory two-dimensional calculations at the seventh radial location
also indicate sensitivity to numerical spatial discretization. Here the
width is the amplitude decay rate $\Gamma=-\operatorname{Im}\rho$ for
the convention $\exp(-i\rho t)$. The reported minimum widths scale
approximately as $h^4$ for a five-point square stencil and $h^8$ for both
a weighted nine-point square stencil and a seven-point triangular stencil.
The leading angular radiation components are $\cos(4\theta)$ for both
square stencils and $\cos(6\theta)$ for the triangular stencil. These
observations are consistent with a perturbative interpretation in which
the leading anisotropic operator correction induces a radiation amplitude
linear in its strength, and the width is quadratic in that amplitude.
The analogous quadratic asymmetry law is established for optical
quasi-BICs by Koshelev et al.~\cite{Koshelev2018}, Eqs.~(2)--(3);
its application here remains a model-specific interpretation.
The reported exponent uncertainties are fit errors, and comparisons of
absorbing layers and shifted eigenvalue solves are sensitivity diagnostics
rather than certified error bounds. The calculations concern one
two-dimensional mode and do not establish three-dimensional or nonlinear
persistence. They motivate explicit stencil and resolution checks before
angular radiation patterns are interpreted as physical structure.

\subsection{Remaining mathematical and experimental tasks}

There are three distinct remaining mathematical tasks. First, a full
operator-theoretic realization would sharpen the operational use of
``embedded'' by specifying the generator, its domain and its essential
spectrum. Second, an independent interval implementation would test the
proof software beyond source review and parameter repetitions. Third,
the additional radial case requires its nonlinear outgoing amplitudes,
while the dipole requires a complete quadrupole certificate to address
the circular polarizations not covered by the conditional monopole
bound in Section~\ref{nl:conditional-monopole}. Neither that
calculation nor the earlier evidence for radial location C0 provides a general
continuum theorem of finite-amplitude localization. A short long-lived
trajectory cannot settle that question.

For the dipole calculation, the finite-problem source--current balance
has been checked without refitting the boundary coefficients. Propagated
bounds for the reconstructed input and full exterior now close the
conditional monopole comparison. The resulting nonzero coefficient
obstructs only the uniformly localized smooth continuation class stated
in Section~\ref{nl:polarization}, under the inherited exact-mode and
analytic hypotheses. An independently implemented enclosure and the
unresolved quadrupole response remain separate next checks.

An experimental analogue would additionally need a dimensional mapping,
an identified platform implementing the two-channel operator, measurable
loss channels and a separation from instrumental linewidth floors. None
is established here. The field model and all displayed parameters remain
dimensionless; no laboratory frequency or length is inferred.

\section{Conclusion}
The strongest results are three specific computer-assisted linear
existence enclosures with analytic exterior treatment, rather than a
claim that every observed width minimum is an exact bound state. A larger
numerical sequence and successive predictions support a thin-wall
phase-matching description, while leaving an infinite-ladder theorem
open. The nonlinear expansion identifies additional radiative channels;
a numerical dipole calculation supports second-harmonic leakage for an
approximate input, and a complete error budget establishes a conditional
nonzero monopole coefficient at the enclosed dipole mode. This does not
cover circular polarization or establish a lifetime. Linear cancellation, approximate
spacing, and nonlinear persistence are separate questions throughout.

\appendix
\section{Prescribed driving and response diagnostics}\label{app:driving}
\paragraph{Linear coupling to the localized mode.}
For radiation of the same scalar field, a conditional projection identity
makes the distinction between scattering and local forcing precise.
In a fixed real $l=1$ harmonic, complexify the doubled reduced perturbation
$\psi=(rz_l,r\overline z_l)^T$ and write an additive source as $S$.
Equation~\eqref{nl:linear} and its conjugate give
\begin{equation}
 \psi_{tt}+2i\omega\sigma_3\psi_t+\mathcal K\psi=S,
 \qquad
 \mathcal K=-\partial_r^2+\frac{2}{r^2}
 +(D-\omega^2)I+C\sigma_1,
 \label{env:time}
\end{equation}
where $D,C$ are those of Eqs.~\eqref{eq:radialA}--\eqref{eq:radialB},
$\sigma_3=\operatorname{diag}(1,-1)$, and
$\sigma_1=\left(\begin{smallmatrix}0&1\\1&0\end{smallmatrix}\right)$.
The physical perturbation is recovered by the conjugate frequency pairing.
For a physical additive source in this real harmonic block,
$S_2=\overline{S_1}$; independent components are used only in the formal
complexification, with the conjugate pairing restored for physical forcing.
Use the radial inner product with measure $dr$ and the regular
self-adjoint domain of $\mathcal K$. With $X=(\psi,\psi_t)^T$, define
\begin{equation}
 \mathcal G=\begin{pmatrix}0&I\\-\mathcal K&-2i\omega\sigma_3\end{pmatrix},
 \qquad
 \mathcal J=\begin{pmatrix}2\omega\sigma_3&-iI\\iI&0\end{pmatrix}.
 \label{env:generator}
\end{equation}
Then $\mathcal G^\dagger\mathcal J+\mathcal J\mathcal G=0$ on that domain.
For a nonzero localized fundamental mode
$H(\rho)u=0$, where $H(\rho)=\mathcal K-\rho^2I-2\omega\rho\sigma_3$,
put $X_u=(u,i\rho u)^T$. If $\rho>|\omega|$, its conserved form is positive:
\begin{equation}
 N=\langle X_u,\mathcal JX_u\rangle
   =-\langle u,H'(\rho)u\rangle
   =2\langle u,(\rho I+\omega\sigma_3)u\rangle>0.
 \label{env:norm}
\end{equation}
Here $\rho$ is the fundamental frequency, not $2\rho$.
Positivity on this mode is neither a linear nor a nonlinear stability
theorem for the full system.
The corresponding left projection $c=\langle X_u,\mathcal JX\rangle/N$
obeys
\begin{equation}
 \dot c=i\rho c-\frac{i}{N}\langle u,S\rangle,
 \label{env:projection}
\end{equation}
provided the radial Green boundary form vanishes. On $[0,R]$ the
additional term is $i[u'^\dagger\psi-u^\dagger\psi']_0^R/N$.
At the origin this follows from regular values $u,\psi=O(r^2)$ and
regular derivatives $u',\psi'=O(r)$. At infinity the localized mode and
its derivative must decay against the controlled scattering field and
its derivative; this includes suitable finite wave packets.

Thus a homogeneous linear solution ($S=0$) with $c(0)=0$ cannot populate
this projection on the fixed background. Spatial separation of an
incoming packet alone does not imply exact orthogonality, because the
localized mode has tails. A prescribed local additive source can instead
have $\langle u,S\rangle\ne0$ and drive the mode. For $c(0)=0$ and a
finite square-integrable projected source, Eq.~\eqref{env:projection} gives
\begin{equation}
 N|c(T)|^2\leq\frac{T}{N}\int_0^T|\langle u,S(t)\rangle|^2\,dt.
 \label{env:pulse}
\end{equation}
This is a mode-projection bound, not a statement about the energy supplied
by an unspecified environment. Although Eq.~\eqref{env:norm} is positive
on this eigenspace, $\mathcal J$ is indefinite on the full phase space.
No globally passive port model, loading rate, nonlinear formation
mechanism, or indefinite maintenance follows. The local source and its
physical coupling must be specified separately; the nonlinear radiation
obstruction in Section~\ref{nl:section} addresses a different question.

% Result, detuning lemma and text independently reviewed; v0.10 stage.
\paragraph{A prescribed local pulse.}
For the exact T2 mode, write $u=(A,B)^T$ with the reduced real components
of Eq.~\eqref{eq:ansatz}, in the exact-mode Jost normalization specified in
Section~\ref{nl:conditional-monopole}. The projection identity gives a
concrete excitation test for this already existing bound mode.
In dimensionless mass-one units, prescribe a source with real angular
profile $Y_{10}$. Choose the radial envelope
$\chi(r)=(r-1)^2(2-r)^2$ on $1\le r\le2$
and zero elsewhere, and temporal envelope $\eta(t)=\sin^2(\pi t)$ on
$0\le t\le1$ and zero otherwise. The physical source in the original field equation is
$J_\phi=\epsilon\chi Y_{10}\eta\exp[i(\omega+\rho)t]$, with real
$\epsilon$, the small strength of the external drive in this linear-response
calculation. It is added to the right-hand side of Eq.~\eqref{nl:model}.
The reduced source in Eq.~\eqref{env:time} is
$S=\epsilon r\chi\eta(e^{i\rho t},e^{-i\rho t})^T$;
its conjugate channel is not independent and must be retained. For the fixed real
radial mode components $A,B$, define
\[
 I_A=\int_1^2 r\chi(r)A(r)\,dr,\qquad
 I_B=\int_1^2 r\chi(r)B(r)\,dr.
\]
The resulting finite-pulse overlap is
\[
L(0)=\frac{I_A}{2}+I_B\int_0^1\eta(t)e^{-2i\rho t}\,dt.
\]
With $N>0$ from Eq.~\eqref{env:norm}, vanishing initial projection and
vanishing Green boundary form, the pulse leaves
$c(1)=-i\epsilon e^{i\rho}L(0)/N$.
Exact integration of the stored nominal polynomials, combined with their
certified componentwise error bounds and the same certified frequency box,
gives an enclosure separated from zero. Conditional on the same exact T2
mode identification and normalization,
a conservative bound is $|L(0)|>1/500$. This certifies nonzero linear mode
projection for nonzero $\epsilon$, vanishing initial projection and the
stated boundary conditions; it does not determine an absolute excitation
energy without the mode norm. Replacing $Y_{10}$ by $Y_{00}$ gives zero
direct dipole projection by angular orthogonality; other modes are not excluded.

% Detuning bound independently checked in DETUNING-REVIEW.txt.
\paragraph{Finite detuning.}
The conclusion need not require an exactly tuned external source. Replacing
its carrier by $\omega+\rho+\delta$ gives
\[
 L(\delta)=I_A\int_0^1\eta(t)e^{i\delta t}\,dt
 +I_B\int_0^1\eta(t)e^{-i(2\rho+\delta)t}\,dt.
\]
Since $\int_0^1t\eta(t)\,dt=1/4$,
\[
 |L(\delta)-L(0)|\le\tfrac14(|I_A|+|I_B|)|\delta|.
\]
The conservative bounds $|I_A|<3/1000$, $|I_B|<6/1000$ therefore imply
$|L(\delta)|>7/8000$ for $|\delta|\le1/2$. This is a sufficient excitation
bound for the prescribed local source, not a linewidth or a decay rate.
Its duration is dimensionless, not a measured physical loading time.
Neither calculation describes formation of the underlying Q-ball,
self-sustained excitation, or coupling to a specified cosmological field.

\paragraph{Canceling one excitation is not stopping the Q-ball.}
The source-projection identity also supplies a simple linear control prediction.
Let $d=e^{-i\rho t}c$, initially zero, and apply two identical physical pulses
of the preceding family, with the same local rotating-frame waveform and
strength $\epsilon$, separated by $\Delta\geq1$.
Their total projected drive is
\begin{equation}
 I_{\mathrm{pair}}=\epsilon L(\delta)
                 \bigl(1+e^{-i\rho\Delta}\bigr).
 \label{env:counterpulse}
\end{equation}
Consequently $\Delta=\pi/\rho>1$ cancels the mode projection generated by the
first pulse, even for a common local detuning $\delta$. Both sideband terms
in $L$ are retained. This construction needs identical pulses in the
rotating frame: in laboratory variables their sources obey
$J_2(t)=e^{i\omega\Delta}J_1(t-\Delta)$.
A sign-reversed second pulse at that same half-period instead doubles the
rotating projected drive. For $|\delta|\leq1/2$, the preceding bound ensures
a nonzero first-pulse projection. Relative to that projection, a timing
error $h$ leaves the remainder $2|\sin(\rho h/2)|\leq\rho|h|$.
These statements are analytic consequences of the fixed linear model.
The following finite-time calculation tests the zero-detuning case only;
it does not numerically establish the entire detuning or timing-error family.
The common unknown modal norm cancels
from the cancellation condition; extinguishing an arbitrary pre-existing
excitation would require its complex amplitude and calibration.
Other modes, continuum waves, radiation already emitted, and nonlinear
backreaction are not canceled by this one condition. Neither zero modal
projection nor zero linear radiation implies that the charged Q-ball itself
has stopped rotating or disappeared.

% Candidate only. No canonical paper edit; figure awaits rendering/review.
\paragraph{A finite-time radial counterpulse test.}
We tested the two-pulse prediction by evolving the linearized complex radial
$\ell=1$ equation on the fixed nominal T2 background, with the conjugate
physical channel retained.
Four initially zero arms used no pulse, one pulse, two equal-sign local
rotating-frame pulses separated by $\Delta=\pi/\rho$, and a sign-reversed
second pulse at the same delay.
The calculation used $R=20$, $T=6$, grid spacings $h=0.04,0.02$ and
fourth-order Runge--Kutta time steps $h/8$, with forcing reported per unit
linear drive amplitude. It is a radial linear PDE calculation, not a
nonlinear three-dimensional evolution.

Both grids passed the predeclared projection-balance and comparison checks.
The equal-sign pair reduced the final target projection below $2\%$ of
the single-pulse reference, while the sign-reversed pair agreed with twice
the reference complex projection within $2\%$.
These are numerical acceptance bounds, not interval-certified errors.
The observed fine-grid ratio was approximately $3.61\times10^{-5}$
($0.0036\%$); the coarse-grid ratio was approximately $1.44\times10^{-4}$.
The balance includes the explicitly computed defect of the stored nominal
mode on the spatial grid; the source-only discrepancy is reported separately.
The global second-component defect norm, which includes the forced cutoff
of the nominal mode at the finite boundary, actually increases on refinement;
we do not claim
global eigenfunction-error convergence. The short-time paired defect and
projection checks concern the evolving, spatially confined response.
Figure~\ref{fig:counterpulse-dynamics} also shows that the radial field
norm need not vanish when the selected projection is suppressed.
At $T=6$ its fine-grid value was $0.0119332$ for the equal-sign pair,
compared with $0.00818567$ for the single pulse.
The mode-subtracted endpoint norm uses the nominal reconstruction
$q-cA-\overline cB$ and is not an orthogonal energy decomposition.
No total energy, charge removal, long-time stability or complete stillness
of the charged background is inferred.

\begin{figure}[htbp]\centering
\includegraphics[width=.98\linewidth]{counterpulse-dynamics.pdf}
\caption{Stored radial linear-PDE histories for the T2 counterpulse test.
Top: target projection magnitude, normalized by the fine-grid final
single-pulse value. Bottom: reduced-field $L^2(dr)$ norms; endpoint symbols
give norms after nominal mode subtraction. Colors identify pulse arms,
solid/dashed lines the two grids. The exactly zero null arm is omitted
from the logarithmic panel. Shading marks pulse support; only paired arms
receive the second pulse. Lines connect recorded times without a new
evolution. The laboratory second pulse includes the carrier-phase correction
of the analytic construction. A small target projection coexists with a
remaining field; these norms are neither radiated energy nor Noether charge.}
\label{fig:counterpulse-dynamics}
\end{figure}

% Candidate for v0.14; not a canonical edit or publication approval.
% RESULT.json SHA256 7c9ff31f20b4ce7ee65e146a13ba1ccdda18ae04c07757e7ad1abf56b4702e88
% RESULT-REVIEW.txt SHA256 4bb05fb4bcaa8c187c7fd850c0d2908283bedf98b3e1d494e5e3676658373e33
\paragraph{Changing the complete pulse shape.}
A subsequent fixed comparison used the same nominal linear radial
$\ell=1$ evolution, $R=20$, $T=6$ and $h=0.04,0.02$.
It compared the annular drive above with a mode-adapted drive, using
single pulses and equal-sign pairs separated by $\pi/\rho$.
The new drive retains both sideband terms of a quintic turn-on construction
and has a spatial cutoff equal to one for $r\leq4$ and zero for $r\geq8$.
The reference instead uses the annulus $1<r<2$ and a $\sin^2$ temporal
window with one rotating frequency. Thus both spatial and temporal
waveforms change; the comparison does not isolate a spatial cause.

Each single pulse was normalized to the same reduced-source norm
\begin{equation}
 B_{\mathrm{src}}=\int\!\int |s(r,t)|^2\,\dd r\,\dd t,
 \qquad B_{\mathrm{src}}\simeq1.35281\times10^{-3}
 \quad(h=0.02).
\end{equation}
Pairs have twice this budget. The integrated budgets agreed with their
prescribed values within $6\times10^{-9}$ relatively on both grids.
This norm is neither mechanical work nor injected energy or charge,
and equal norms do not imply equal actuator support or hardware cost.
For the single pulse, the fine-grid diagnostic $|d(T)|^2/B_{\mathrm{src}}$
is about $32$ times larger for the shaped drive
(Table~\ref{tab:source-shaping}); this is not an energy-efficiency ratio
or an optimality result.

\begin{table}[htbp]\centering\small
\caption{Fixed source-shaping comparison at $T=6$, $h=0.02$.
Here $p=\partial_tq$, and ``remainder'' means subtraction of the nominal
mode reconstruction. Field norms are $L^2(\dd r)$ norms, not energies;
the pair rows give the total remaining field and velocity norms.}
\label{tab:source-shaping}
\begin{tabular}{lrr}\toprule
Diagnostic&Annular drive&Shaped drive\\\midrule
Single-pulse $|d|^2/B_{\mathrm{src}}$&$0.001214$&$0.03872$\\
Single-pulse $q$ remainder norm&$0.007992$&$0.002596$\\
Single-pulse $p$ remainder norm&$0.01502$&$0.004954$\\
Pair $\|q\|$&$0.01193$&$0.001055$\\
Pair $\|p\|$&$0.02153$&$0.001930$\\\bottomrule
\end{tabular}
\end{table}

The shaped pair leaves approximately $8.8\%$ and $9.0\%$ of the
reference pair's field and velocity norms, respectively; both remain
nonzero. Nominal mode subtraction is not an orthogonal energy
partition. Both grids passed the fixed comparison and defect-balance
checks, but the global second-component nominal eigenmode residual norm increases
from approximately $0.529$ to $1.495$ on refinement. The successful
paired-defect balance therefore does not establish global eigenfunction
convergence. These finite-box, short-time results show reduced residual
motion for the specified composite drive; they establish neither full
stillness of the field nor removal of the charged background, nonlinear
stability or a lifetime. This original two-source comparison did not
separate changes of spatial and temporal design. The subsequently
specified cross-comparison below addresses that limitation within a
fixed actuator family.

% Reviewed stored-profile attribution, incorporated in draft 0.17.
% RESULT 03e1c3129984f7b96713f901e7e7767b5f32d4bb6b33afa380178d34dc8ffbb7
% RESULT-REVIEW 3b665e564d1471fcfbf80cb42d118615c011bfb7b4a2cffea75522ca27619b9c
A subsequent stored-profile audit identifies the origin of the growing
second-component defect in this box. The exported reference mode has
$v(20)=0.004228503237$, whereas the evolution operator imposes a zero
Dirichlet value. For the three-point stencil this discrepancy adds
$b_j=\delta_{j,n-1}v(20)/h^2$ to the continued-ghost residual, giving
$\|b\|_h^2=|v(20)|^2/h^3$. Its squared norm is $0.2793787442$ and
$2.2350299533$ on the two grids; the continued residual squared norms
are $1.52\times10^{-8}$ and $9.50\times10^{-10}$, with signed cross
terms $-9.74\times10^{-10}$ and $-4.84\times10^{-10}$.
The boundary discrepancy therefore dominates this component's growing
residual. It does not dominate the first component. This is an algebraic
attribution, not a modification of the evolution boundary condition or a
new eigenfunction calculation. The original global defect remains valid
for the original operator. Two decreasing interior residuals neither
certify the exported mode nor bound the integrated influence of the
boundary discrepancy on the pulse response.

\begin{figure}[htbp]\centering
\includegraphics[width=.98\linewidth]{source-shaping.pdf}
\caption{Endpoint diagnostics from the fixed spatiotemporal pulse comparison.
Both spatial support and temporal waveform change. The common source norm
is not energy or work. The shaped source increases the selected projection
per source norm and reduces the residual response in this comparison;
the complete field does not vanish. Bars show the stored fine-grid result,
not a new evolution or a certified continuum limit.}
\label{fig:source-shaping}
\end{figure}

% Factorial RESULT477978704a19d5b2f2b474c41cf5917367335c521902e0f712303f7a563ec263
% Non-author RESULT-REVIEW66fe64890bbbc3568c19a33610968cd7360f68462ea8f7fb7b467b900dd02a28
\paragraph{A fixed factorial actuator comparison.}
The new comparison uses two spatial pairs, $A=(w,0)$ and
$B=(\zeta u,\zeta v)$, where $w$ is the preceding annulus profile,
$\zeta$ the cutoff defined above (one below $r=4$ and zero above $r=8$),
and $u,v$ the stored nominal sidebands.
Two temporal pairs are prescribed on $0<t<1$:
\begin{align}
 T_0&=\bigl(\sin^2(\pi t)e^{i\rho t},
                  \sin^2(\pi t)e^{-i\rho t}\bigr),\\
 T_1&=\bigl([g''+2i(\rho+\omega)g']e^{i\rho t},
                  [g''-2i(\rho-\omega)g']e^{-i\rho t}\bigr),
 \qquad g=10t^3-15t^4+6t^5.
\end{align}
Both temporal pairs vanish outside $0<t<1$; $g$ is extended constantly
as zero before the pulse and one afterward.
Each complex source is the component sum
$s_{ij}=\mathcal N_{ij}(R_{i1}T_{j1}+R_{i2}T_{j2})$,
where $R_i=A$ or $B$. There is no conjugation in this sum; the second
physical equation receives $s_{ij}^*$.
The positive factor $\mathcal N_{ij}$ fixes the same single-pulse
$B_{\mathrm{src}}$ \emph{after} summation, retaining component interference.
The two old corners $A/T_0$ and $B/T_1$ are implementation controls;
$A/T_1$ and $B/T_0$ are the new crossed drives. Each uses both a single
pulse and an equal-sign pair at the unchanged delay $\pi/\rho$.
The nominal operator, $R=20$, $T=6$, $h=0.04,0.02$ and time step $h/8$
remain fixed. Both grids passed the predeclared numerical checks.

\begin{table}[htbp]\centering\small
\caption{Four fixed normalized actuator designs at $T=6$, $h=0.02$.
The projection diagnostic is for a single pulse; the final two columns
are total remaining pair norms, not mode-subtracted norms or energies.}
\label{tab:source-factorial}
\begin{tabular}{lrrr}\toprule
Design&$|d_{\mathrm{single}}|^2/B_{\mathrm{src}}$&Pair $\|q\|$&Pair $\|p\|$\\\midrule
$A/T_0$&$0.001214$&$0.01193$&$0.02153$\\
$A/T_1$&$0.0004100$&$0.006390$&$0.01426$\\
$B/T_0$&$0.1586$&$0.02024$&$0.005692$\\
$B/T_1$&$0.03872$&$0.001055$&$0.001930$\\\bottomrule
\end{tabular}
\end{table}

Within this family, $B/T_0$ gives the largest single-pulse projection
diagnostic, whereas $B/T_1$ leaves the smallest pair norms in both $q$
and $p$. Strong projected excitation and small residual field are thus
different objectives. The spatial-pair change $A\to B$ increases the
pair $q$ norm at $T_0$ but decreases it at $T_1$: its effect depends on
the temporal partner. These are conditional design comparisons, not
universal additive spatial and temporal effects. In particular $A$
omits a spatial component, $T_1$ changes sideband weights and phase
structure, and normalization depends on both factors. Equal source norm
still implies neither equal work nor energy efficiency. All four pairs
suppress their selected projection under the fixed criterion, while their
full responses differ and remain nonzero. The nominal-mode defect and
finite-box limitations above remain. That fixed-phase comparison did not
test phase-error robustness; it established neither exact preparation,
complete stillness nor a lifetime.

\begin{figure}[htbp]\centering
\includegraphics[width=.98\linewidth]{factorial-comparison.pdf}
\caption{Stored endpoint diagnostics for the fixed $2\times2$ actuator
family. The largest single-pulse projection per source norm occurs at
$B/T_0$, while $B/T_1$ gives the smallest remaining pair field and
velocity norms. The source is normalized after adding its two components;
these comparisons are not energy efficiencies or a universal separation
of independent spatial and temporal mechanisms. No additional evolution
or phase-error test is represented by this figure.}
\label{fig:source-factorial}
\end{figure}

% Phase RESULTf492a34c52a608a36d3d6054c1342cc6bf3dd963021e3576d5d6447423a25301
% Reviewed ea5b091239d228e8da5ac68c60108c238c3052a2ee4272f9012366e61459ec37
% Endpoint SUMMARY3329b5213520fe157b71c0a79d530a36e03ce9a644fd13d48cc71031e0c53aae
\paragraph{A constant relative pulse phase.}
A further bounded test changes only the additional phase $\theta$ of the
second pulse, keeping the delay and the carrier-phase correction fixed.
The coupling to $q^*$ makes the nominal Bogoliubov evolution real-linear,
not complex-linear: the response to $i s$ need not equal $i q$.
For each of the four fixed designs, one new delayed-$i s$ response
therefore completes the real basis. At the same endpoint,
\begin{equation}
 q_\theta=q_{\mathrm{single}}
   +\cos\theta\,(q_{\mathrm{pair},0}-q_{\mathrm{single}})
   +\sin\theta\,q_{\mathrm{delayed},90}.
 \label{eq:phase-real-basis}
\end{equation}
The same identity holds for $p$ and the real-linear projection $d$.
Four directly evolved $\theta=\pi/4$ pairs check this reconstruction:
the largest reported relative discrepancy in $q,p,d$ across both grids
was $1.14\times10^{-13}$. This tests a known linear identity, not a new
physical effect. Other constant phases below are evaluations of stored
endfields, not additional time evolutions. The setting remains
$R=20$, $T=6$, $h=0.04,0.02$.

\begin{table}[htbp]\centering\small
\caption{Selected constant-phase endpoints for $B/T_1$ at $T=6$,
$h=0.02$. Norms are total remaining radial fields, not energy.
The phase values were fixed before the endpoint analysis.}
\label{tab:phase-selected}
\begin{tabular}{rrr}\toprule
$\theta$&$\|q\|$&$\|p\|$\\\midrule
$0^\circ$&$0.001055$&$0.001930$\\
$-5^\circ$&$0.001062$&$0.002068$\\
$+5^\circ$&$0.001423$&$0.002462$\\
$-15^\circ$&$0.001986$&$0.003720$\\
$+15^\circ$&$0.002562$&$0.004413$\\
$45^\circ$&$0.006301$&$0.01125$\\\bottomrule
\end{tabular}
\end{table}
These values need not be symmetric about zero phase. Cancellation of the
selected coefficient and minimization of the full field norm are different
conditions; no new minimum search was performed.

For a uniformly distributed phase that is \emph{constant within each
pulse pair}, write any endpoint field as $a+b\cos\theta+c\sin\theta$.
Its root mean square norm is exactly
\begin{equation}
 \sqrt{\mathbb E_\theta\|a+b\cos\theta+c\sin\theta\|^2}
 =\sqrt{\|a\|^2+\tfrac12(\|b\|^2+\|c\|^2)}.
\end{equation}
The Gram-matrix evaluation was also checked against eight equally spaced
phases. Table~\ref{tab:phase-rms} keeps all four designs in the comparison.
\begin{table}[htbp]\centering\small
\caption{Uniform constant-phase RMS endpoint norms, $T=6$, $h=0.02$.
These are square roots of mean squared norms, not mean norms.}
\label{tab:phase-rms}
\begin{tabular}{lrr}\toprule
Design&$\sqrt{\mathbb E_\theta\|q\|^2}$&$\sqrt{\mathbb E_\theta\|p\|^2}$\\\midrule
$A/T_0$&$0.01128$&$0.02168$\\
$A/T_1$&$0.007123$&$0.01447$\\
$B/T_0$&$0.02552$&$0.04088$\\
$B/T_1$&$0.01106$&$0.01954$\\\bottomrule
\end{tabular}
\end{table}
Although $B/T_1$ leaves smaller pair norms than $A/T_1$ at zero phase,
its two RMS norms are larger in this ensemble. Its fixed-phase advantage
is thus conditional on phase preparation.
The comparison fixes source norm, not the previously excited target
amplitude: $A/T_1$ has a much smaller single-pulse projection diagnostic
($0.0004100$ versus $0.03872$ for $B/T_1$). Its lower RMS therefore does
not establish a better storage or erasure protocol at equal preparation.
An RMS ranking is not necessarily
an ordering of mean norms. This ensemble is neither a time average nor a
simulation of fluctuating environmental noise. The constant phase does
not change the disjoint pulses' source budget $2B_{\mathrm{src}}$, but that
budget remains distinct from work or energy. The global nominal eigenmode
residual limitation remains unchanged; neither exact BIC preparation,
complete stillness nor nonlinear or long-time robustness follows.

\begin{figure}[htbp]\centering
\includegraphics[width=.98\linewidth]{phase-response.pdf}
\caption{Constant-phase endpoint reconstruction for all four fixed
actuator designs, from their stored real-linear response bases.
The two grids compare nominal finite-box observables; plotted phase
samples are not independent time evolutions. The additional phase is
constant within a pulse pair, not a time-dependent noise source.
Field norms and uniform-phase RMS values do not measure energy cost.}
\label{fig:phase-response}
\end{figure}

\paragraph{Quadratic endpoint energy and the rotating generator.}
A further analysis uses the stored single- and pair-pulse endpoints,
without repeating the time evolution. With $p=q_t$ and the same
$h$-weighted inner product and Dirichlet operator as the evolution,
define
\begin{align}
 H_2&=\|p\|_h^2+\langle q,Kq\rangle_h
       +\operatorname{Re}\,h\sum_j C_j q_j^2,
       \label{eq:endpoint-generator}\\
 Q_2&=2\operatorname{Im}\langle q,p\rangle_h
       +2\omega\|q\|_h^2,\qquad E_2=H_2+\omega Q_2.
       \label{eq:endpoint-energy}
\end{align}
Here $K=-\Delta_h+2/r^2+1-4f^2+(9/2)f^4-\omega^2$ and
$C=-2f^2+3f^4$. There is no factor $1/2$: $H_2$ is the coefficient
of second order in $E-\omega Q$ for the stated action convention,
using one physical complex field rather than counting its conjugate twice.
This rotating quadratic generator need not be positive.

\begin{table}[htbp]\centering
\caption{Stored endpoint coefficients at $T=6$, $h=0.02$.
All entries are in units of $10^{-4}$, rounded to four significant digits.
Within a fixed pulse count the designs have equal source-L2 budgets;
a pair has twice the single-pulse budget.}
\label{tab:endpoint-energy}
\begin{tabular}{lrrrr}\toprule
Design&$H_2$ single&$H_2$ pair&$E_2$ single&$E_2$ pair\\\midrule
$A/T_0$&4.782&9.315&5.958&11.51\\
$A/T_1$&2.668&5.180&3.614&6.910\\
$B/T_0$&8.917&1.382&6.467&8.608\\
$B/T_1$&2.129&0.05203&1.201&0.03759\\\bottomrule
\end{tabular}
\end{table}
The distinction matters even for the direction of change:
$B/T_0$ decreases $H_2$ while increasing $E_2$ from single to pair.
Its $Q_2$ changes from approximately $-2.821\times10^{-4}$ to
$8.319\times10^{-4}$. For $B/T_1$ both endpoint coefficients decrease;
for the two $A$ designs both increase. These directions also occur on
$h=0.04$, but two-grid differences are not validated error bounds.

Equal source-L2 budgets do not imply equal work, and the single-to-pair
comparison itself is not a same-budget comparison. For the fixed
semidiscrete model, $\dot H_2=2\operatorname{Re}\langle p,S\rangle_h$.
Thus $H_2(T)$ at zero initial data predicts net rotating-generator work
by an identity; that endpoint-only analysis did not measure an independent
time-integrated power balance. The subsequent test below adds this observable.
The quadratic laboratory source power additionally contains
$2\omega\operatorname{Im}\langle q,S\rangle_h$, while $\dot E_2$
also contains the background-exchange term
$2\omega\operatorname{Im}\,h\sum_j C_jq_j^2$.
In particular, $E_2$ is a coefficient on the linear ansatz, not the complete
laboratory energy change including a second-order background response.
No positive or gross actuator work, energy efficiency, or full shutdown
of the Q-ball is established by Table~\ref{tab:endpoint-energy}.
\paragraph{Time-integrated source work in the nominal linear model.}
A subsequent bounded run retains the same eight single/pair arms,
$R=20$, $T=6$, $h=0.04,0.02$ and $\Delta t=h/8$, and adds power
accumulators at the four Runge--Kutta stages. With the conventions of
Eqs.~\eqref{eq:endpoint-generator}--\eqref{eq:endpoint-energy}, these
integrate
\begin{align}
 P_H&=2\operatorname{Re}\langle p,S\rangle_h,&
 J_Q&=2\operatorname{Im}\langle q,S\rangle_h,&
 X_Q&=2\operatorname{Im}\,h\sum_jC_jq_j^2,\\
 P_{E,\mathrm{src}}&=P_H+\omega J_Q,&
 P_{E,\mathrm{ex}}&=\omega X_Q.
\end{align}
The independently evaluated state coefficients satisfy the numerical
balances $\Delta H_2=\int P_H\,\dd t$,
$\Delta Q_2=\int(J_Q+X_Q)\,\dd t$, and
$\Delta E_2=\int(P_{E,\mathrm{src}}+P_{E,\mathrm{ex}})\,\dd t$.
Unlike the preceding endpoint-only analysis, these integrals use the
actual time-dependent fields and source at the integration stages.
The exchange term remains active after the source ends.

The first pulse is identical in the single and pair arms; the first and
second pulse windows are disjoint. Their trajectories therefore agree before
the second pulse. Subtracting their independently integrated total source
works gives the second pulse's net work on the already excited field:
\begin{equation}
 W_{H,2}=W_{H,\mathrm{pair}}-W_{H,\mathrm{single}},\qquad
 W_{E,\mathrm{src},2}=W_{E,\mathrm{src},\mathrm{pair}}
                       -W_{E,\mathrm{src},\mathrm{single}}.
\end{equation}
This is not the work of an isolated second pulse applied to a fresh
zero state, nor merely a relabeling of endpoint energy differences.

\begin{table}[htbp]\centering
\caption{Net second-pulse source work inferred from time-integrated
pair-minus-single works at $h=0.02$. Units are $10^{-4}$, as in
Table~\ref{tab:endpoint-energy}; entries are rounded to four significant
digits. Negative values denote net extraction in the specified quadratic
source-work convention, not efficiency of a physical actuator.}
\label{tab:second-pulse-work}
\begin{tabular}{lrr}\toprule
Design&$W_{H,2}$&$W_{E,\mathrm{src},2}$\\\midrule
$A/T_0$&$4.533$&$4.207$\\
$A/T_1$&$2.512$&$3.602$\\
$B/T_0$&$-7.535$&$5.939$\\
$B/T_1$&$-2.077$&$-1.122$\\\bottomrule
\end{tabular}
\end{table}
Thus the second $B/T_1$ pulse is a net source-work sink in both
conventions. In contrast, $B/T_0$ removes rotating-generator work while
its laboratory source-work coefficient is positive. The change in $E_2$
also includes the difference of integrated exchange terms; it is not
identical to $W_{E,\mathrm{src},2}$.
Positive and negative parts of the total spatially integrated power
were separately accumulated, so a small net value need not mean little
positive input or little extraction. This splitting applies after spatial
summation, rather than to positive and negative local power densities.

Both grids pass the predeclared $10^{-6}$ normalized balance tolerance,
and the new endfields agree with the stored reference endpoints within
the fixed $10^{-8}$ single-response-normalized tolerance. The reported
scaled two-grid differences of work contributions are below
$4.4\times10^{-4}$, within the prescribed $0.03$ consistency limit.
These are discrete numerical checks, not rigorous continuum error bounds;
the positive/negative power split also need not inherit a smooth fourth-order
convergence rate. The original nominal-mode boundary defect is unchanged.
Equal source-L2 budgets are still not equal work. Finally, the laboratory
source-work coefficient and $E_2$ belong to the frozen linear ansatz:
second-order background response, full nonlinear laboratory energy,
actuator losses and conversion efficiency are not determined. No complete
shutdown or long-time stability follows from this short-window balance.
\paragraph{Phase dependence of the second-pulse source work.}
For the same four source designs, we vary the constant phase $\theta$ of
the delayed pulse while keeping its source norm fixed. Real linearity of
the frozen evolution gives the quadratic source-work coefficient
\begin{equation}
 W_{E,\mathrm{src},2}(\theta)
 =c_0+c_1\cos\theta+s_1\sin\theta
       +c_2\cos(2\theta)+s_2\sin(2\theta).
 \label{eq:phase-work-five}
\end{equation}
Writing the measured quadrature coefficients as $(A,B,D,E,F)$, these are
$(c_0,c_1,s_1,c_2,s_2)=((D+E)/2,A,B,(D-E)/2,F/2)$.
The five coefficients are integrated from the first-pulse response and two
delayed quadrature responses, with a direct $45^\circ$ two-pulse control,
through $T=3$ after both source windows end ($\Delta+1<3$), using
$h=0.04,0.02$ and $\Delta t=h/8$. This is not a renewed $T=6$ energy history.
Thus the plotted curves are postprocessing of measured coefficients, not
thousands of independent evolution runs.

On both nominal grids, only B/T1 has a negative-work interval:
approximately $-66.53^\circ<\theta<57.03^\circ$.
For the fine grid its second-pulse net source work is
$-1.1224\times10^{-4}$ at $0^\circ$ and
$-3.9278\times10^{-5}$ in the direct $45^\circ$ control.
This broad nominal interval concerns signed net source work; it does not
assert equally small residual fields throughout the interval.
In particular, the previously reported residual field norms increase at
$45^\circ$ even though this work remains negative.
All four uniform \emph{constant-phase} means $c_0$ are positive;
for B/T1 the fine-grid mean is $1.1334\times10^{-4}$.
This phase ensemble is not a model of time-dependent phase noise.

Polynomial sign isolation treats the measured floating-point coefficients
as exact rational dyads. Its exactness applies to those nominal polynomials,
not to their discretization errors or to the continuum or nonlinear field
equations. The result concerns quadratic source work in the frozen linear
ansatz, with background exchange distinguished separately; it is neither a
full nonlinear laboratory-energy balance nor an actuator-efficiency claim.

\begin{figure}[htbp]
\centering
\includegraphics[width=\linewidth]{phase-work.pdf}
\caption{Second-pulse source work versus constant phase. Left: all four
designs; solid curves use $h=0.02$, dashed curves $h=0.04$.
Right: B/T1 with nominal fine-grid zero boundaries (dotted), its uniform
constant-phase mean (dash-dotted), and direct $45^\circ$ controls (symbols).
Work is displayed in units of $10^{-4}$. Shading denotes the nominal
negative-work window, not certified physical tolerances.}
\label{fig:phase-work}
\end{figure}

\clearpage
\paragraph{One prescribed temporal phase drift.}
The constant-phase ensemble above does not determine the response to a
phase varying during a pulse. As a separate finite comparison, we replaced
the same delayed B/T1 source $S_d(t)$ by
$S_{\delta,\theta_0}(t)=\exp[\ii\theta_0+\ii\delta(t-\Delta)]S_d(t)$,
with the single preselected value $\delta=0.1$ against $\delta=0$.
The multiplier acts on the entire complex source; the original envelope,
carrier construction and delay are unchanged. It preserves the source
magnitude pointwise and its integrated $L^2$ budget. For each detuning we
independently evolved the responses to $S_{\delta,0}$ and
$\ii S_{\delta,0}$ through $T=3$. Real linearity gives the response at
$\theta_0$ as the cosine--sine combination of these two responses;
the second is not assumed to be $\ii$ times the first.

Averaging uniformly over the \emph{constant initial} phase $\theta_0$
eliminates both the mixed quadrature term and the work cross term with
any deterministic first-pulse field. Consequently the mean second-pulse
laboratory source work is one half of the sum of the two isolated
quadrature works. No first-pulse evolution was repeated; this reduction
applies to the mean, not to phase-specific two-pulse work. With
$(h,\Delta t)=(0.04,0.005)$ and $(0.02,0.0025)$, both means are positive
on both discretizations. On the finer grid the mean increases from
$1.133373236\times10^{-4}$ to $1.248632464\times10^{-4}$,
a difference of $1.152592277\times10^{-5}$ (about $10.2\%$).

The increase exceeds the predeclared diagnostic comparison guard
$7.94\times10^{-11}$, which combines the change between discretizations,
unshifted-work reproduction errors and energy-balance residuals.
This guard is not a certified error bound; space and time resolution change
together. The result establishes sensitivity of this nominal mean source
work to one prescribed phase drift at equal source norm. It does not give
a phase-noise or thermal model, an efficiency gain, or a claim about all
detunings. The frozen-profile approximation, artificial $R=20$ boundary
and omitted nonlinear backreaction remain unchanged.
\paragraph{Local direction and one mirror detuning.}
A follow-up fixed $d=0.1$ and added only the mirror source at $-d$ and
the parameter tangent at zero. Writing $z=\partial_\delta q|_0$, the
exact differentiated finite-grid equation has the same real-linear
operator and forcing $\ii(t-\Delta)S_d$. Its numerical integration retains
both derivatives of the work integrand, from the response and from the
source, for each independently evolved quadrature. The zero-detuning
basis trajectories are needed for these products; the first pulse is
not repeated, and the $+d$ result is reused. For the uniform initial-phase
mean $m(\delta)$, the finer grid gives $m'(0)\simeq1.1006065\times10^{-4}$
in model units. Both grids resolve a positive slope with the prescribed
empirical guard. The three fine-grid means remain positive:
\begin{center}
\begin{tabular}{rr}
\toprule
$\delta$ & Mean laboratory source work $m(\delta)$\\
\midrule
$-0.1$ & $1.0285198\times10^{-4}$\\
$0$ & $1.1333732\times10^{-4}$\\
$+0.1$ & $1.2486325\times10^{-4}$\\
\bottomrule
\end{tabular}
\end{center}

Define the finite odd and even changes by
$O_d=[m(d)-m(-d)]/2$ and
$E_d=[m(d)+m(-d)-2m(0)]/2$.
Their fine-grid values are approximately $1.1005634\times10^{-5}$ and
$5.20289\times10^{-7}$; the residual
$R_d=O_d-dm'(0)\simeq-4.31\times10^{-10}$ is small but not an exact
identity or a certified Taylor remainder. The resolved positive odd
change agrees with the local direction. This is a directional response
of the fixed nominal actuator, not an identification of rotation as its
cause: a free complex accelerator $q_{tt}=\exp[\ii(\kappa+\delta)t]$ already
admits a nonzero detuning slope. More generally, uniform initial-phase
averaging selects the normal causal response, whose work derivative is
a lag-weighted imaginary response kernel; the carrier and field response
enter it jointly. The two-grid guards are diagnostic, and the inherited
profile, boundary and linearization limitations persist. No thermal,
new-memory-mediator or efficiency conclusion follows.

\paragraph{Prescribed finite phase coherence.}
A fluctuating drive phase can be specified without identifying it with
thermal noise. Let the finite-duration source be
$S_\theta(t)=e^{i\theta(t)}S_d(t)$, where the initial phase is uniform and
independent and $\theta(t)-\theta(\Delta)=\sqrt{2D}\,B_{t-\Delta}$ during
the pulse, with $B$ a standard real Wiener process independent of that initial phase. Then
$\mathbb E e^{i[\theta(t)-\theta(s)]}=e^{-D|t-s|}$ and the anomalous
source correlation vanishes. Every realization retains the source's
$L^2$ budget. This does not fix the work delivered to the field.

For the fixed finite real-linear response used above, let $m(\delta)$ denote
the phase-averaged source work when the \emph{entire} prescribed source
is multiplied by $e^{i\delta(t-\Delta)}$. Its phase-diffusion mean obeys
\[
 m_D=\int_{-\infty}^{\infty}
 \frac{D}{\pi(D^2+\delta^2)}\,m(\delta)\,d\delta,
 \qquad D>0.
\]
This follows from the characteristic function of the Cauchy distribution
and the quadratic response kernel. At finite time the finite-dimensional
propagator is bounded; compact source support and finite $L^2$ norm give
absolute integrability and justify interchanging the integrals.
The same weighting convolves the finite-pulse spectral power and preserves
its integrated weight. Since the weight is even, the odd detuning response
cancels. The positive local slope found above therefore does not determine
whether phase diffusion increases or decreases mean work. Nor does a single
symmetric detuning pair settle this question.

For computation of this mean work, the lag factor $e^{-D(t-s)}$ can instead
be included in the auxiliary generator $A-DI$, evaluated for two genuine
source quadratures. With state coordinates $(q,p)$, both equations acquire
$-D$ terms; the work output remains $p+i\omega q$, although now
$\dot q=p-Dq$. These auxiliary states are not stochastic field trajectories,
and their endpoint energies are not ensemble energies. This reduction
concerns the specified work observable and does not insert physical friction
into the original field equations.

A Cauchy-distributed constant detuning and a Wiener phase process share the
relevant two-point statistics but not their trajectory laws. Their quadratic
response means agree; the auxiliary damped endpoint is a different object.
Ideal Wiener phase also has no finite ultraviolet cutoff. Thus this identity
is a finite-model interpretation, not a continuum certificate, a temperature
assignment, or a derivation of an environmental bath. Physical bandwidth and
coupling require separate specification before experimental use.
\paragraph{One finite-coherence comparison.}
We fixed $D=1$, corresponding to one pulse duration for the exponential
phase-correlation time, and compared it with $D=0$. The two genuine
quadratures of each auxiliary response were evolved on the same two
discretizations as the detuning test. The finer grid gives mean source
works $m_0=1.133373236\times10^{-4}$ and
$m_1=2.173314375\times10^{-4}$. Their difference,
$1.039941139\times10^{-4}$, is positive on both grids and exceeds the
predeclared diagnostic guard $6.63\times10^{-9}$. The undiffused adapter
reproduces the earlier work values, and the modified auxiliary balances
include $-2D$ times each quadratic form. Free and restoring-oscillator
controls test both the work output and the placement of damping in both
auxiliary coordinates.

This is an increase in mean work supplied by the prescribed external
source at fixed source $L^2$ budget. It is not a gain in energetic efficiency
or a demonstration of self-sustained oscillations. The guard combines
grid differences and implementation residuals; it is not a certified
error bound, and space and time resolution change together. The nominal
profile, finite boundary and linear-response limitations persist. Only
one nonzero diffusion constant is compared: no noise optimum, temperature,
nonlinear formation or lifetime conclusion is established.
The same source-output reduction with output $p$ gives an additional
observable. In the undamped physical finite model, the rotating quadratic
generator satisfies $\dot H=2\operatorname{Re}\langle p,S_\theta\rangle$.
For zero initial perturbation, its ensemble endpoint mean therefore equals
the half-sum of the two auxiliary integrals of this \emph{original} source
rate. It does not equal the auxiliary endpoint $H$. Extracting these
already stored integrals gives, on the finer grid,
$\mathbb E H(T)=2.09020168\times10^{-4}$ for $D=0$ and
$2.93069596\times10^{-4}$ for $D=1$. Both grids show this approximately
$40\%$ increase; this is a descriptive follow-up without a new certified
comparison bound or field integration.

This rotating quadratic generator need not be positive, and its global
value does not measure energy localized near the Q-ball. Full laboratory
endpoint energy additionally requires the physical ensemble charge,
including the anomalous state-exchange contribution. The latter is not
given by the exchange evaluated on the damped auxiliary states. Thus
neither the source-work increase nor this rotating-generator mean establishes
laboratory energy retention or improved confinement.

\paragraph{Correlated instantaneous frequency.}
A separate prescribed drive model replaces Wiener phase increments by
$\theta(t)=\theta_0+\int_\Delta^t\xi(s)\,ds$, with an independent uniform
initial angle and stationary zero-mean Ornstein--Uhlenbeck frequency
$\mathbb E\xi(t)\xi(s)=(D/\tau_c)e^{-|t-s|/\tau_c}$.
Its phase correlation at lag $u\ge0$ is
\[
 C(u)=\exp\{-D[u-\tau_c(1-e^{-u/\tau_c})]\}.
\]
This retains the pathwise source norm and fixes the long-time phase
diffusion constant $D$; the instantaneous frequency variance is
$D/\tau_c$. The Wiener kernel is recovered as $\tau_c\to0$ at fixed $D$.
Finite $\tau_c$ removes its linear short-lag cusp, but does not impose a
hard bandwidth cutoff or define a thermal bath. Pointwise larger phase
correlation need not give larger work against an oscillatory response.

\paragraph{Small correlation-time limit.}
For a fixed finite spatial operator and fixed $D\geq0$, let $\tau=\tau_c$ and write the exact mean work as
$m_\tau=\int_0^T W(u)C_\tau(u)\,du$, with real lag kernel $W$ and
$m_D=\int_0^T W(u)e^{-Du}\,du$. If $|W(u)|\leq M$, then
\[
 |m_\tau-e^{D\tau}m_D|\leq
 \frac{e^{D\tau}DM\tau^2}{1+D\tau},\qquad m'_+(0)=Dm_D.
\]
If $W$ is continuous at zero, the present source-work observable has
$W(0)=2B_{\mathrm{src}}$, where $B_{\mathrm{src}}=\int_0^T\|S(t)\|^2dt$ is the exact
nominal source budget. Consequently,
$m_\tau=m_D+Dm_D\tau+(D^2m_D/2-2DB_{\mathrm{src}})\tau^2+o(\tau^2)$.
These local asymptotics neither determine the comparison at $\tau=0.1$
nor imply global monotonicity, a thermal interpretation, or a continuum error bound.

For the fixed choice $D=1$, $\tau_c=1/10$, put $a=1/10$.
The mean source work has the signed representation
$m_{\mathrm{OU}}=e^a\sum_{n\ge0}(-a)^n m_{1+10n}/n!$, where $m_\lambda$
is the two-quadrature work obtained from the auxiliary generator
$A-\lambda I$. This is not a positive mixture of dissipative states.
The selected truncation retains $n=0,\ldots,8$. If
$B\ge2\int_{t>s}|K_{\mathrm{work}}(t,s)|\,dt\,ds$, its analytical work-tail
bound is
\[
 \epsilon_{\mathrm{tail}}\le
 B\,\frac{10}{9}\,\frac{(1/10)^9}{9!}\,\frac{1}{1-1/100}.
\]
The finite-model bound uses a positive reference energy norm,
a coefficient-envelope check and an exact-rational source-norm bound
for frozen nominal binary coefficients. This bounds the omitted series,
not floating exponential evaluation, signed-sum roundoff or RK integration.
Those errors and the combined space/time sensitivity require separate
assessment; auxiliary endpoint energies are not ensemble energies.
The first finite-correlation implementation passed its scalar controls but
failed the preselected modified quadratic-balance tolerance on the coarser
field grid. That implementation therefore yielded no accepted finite-correlation
field result, and its finer grid was not run. The analytical series-tail estimate
did not repair this numerical acceptance failure, and the original outcome
and thresholds remain unchanged. A separate calculation with a different
integrator and a prospectively specified temporal-sensitivity comparison is
reported below.

\paragraph{A separate finite-correlation pulse comparison.}
A separately specified two-stage Gauss calculation retained the prescribed
source, its normalization, $D=1$, $\tau_c=0.1$ and the signed nine-term expansion.
It split integration steps at the two known source-window joins and used
both genuine quadratures for every auxiliary rate, plus a zero-source
control. Three fixed integrations compared $(h,\Delta t)=(0.04,0.005)$,
$(0.04,0.0025)$ and $(0.02,0.0025)$ through $T=3$ at $R=20$.
All 108 recorded acceptance checks passed; this comparison was made against
the previously stored Wiener-source work on the corresponding spatial grid,
not against an auxiliary endpoint energy or a newly repeated Wiener run.

On the finer spatial grid, the truncated OU mean source work was
$2.1932622828\times10^{-4}$, compared with $2.1733143747\times10^{-4}$
for the stored Wiener result. The difference, $1.9947908110\times10^{-6}$
(approximately $0.918\%$), exceeded the predeclared diagnostic guard
$1.2658579115\times10^{-8}$ and had the same positive sign on the coarser grid.
The guard includes the analytical series-tail allowance and measured temporal,
spatial, adapter and accounting contributions; it is not a certified total-error
bound. The fixed coarse-grid time comparison does not directly measure fine-grid
temporal error. This is a bounded increase in mean source work for one prescribed
drive in the frozen finite-domain model, not evidence of retained energy,
heating, formation, improved efficiency or experimental realizability.

\begin{figure}[htbp]
\centering
\includegraphics[width=.98\linewidth]{mean-injected-work.pdf}
\caption{Cumulative mean source work for prescribed OU and Wiener phase
drives on the $h=0.02$ grid (left), and their difference (right).
The OU curve is the signed nine-term mean of the two genuine quadrature
works per auxiliary rate; the Wiener curve uses the previously stored
calculation. Both use integrated original source work, not auxiliary
endpoint or ensemble energy. Straight lines join 31 stored times, with no
additional field integration for the figure. The endpoint marker spans
the predeclared diagnostic guard; it is neither a confidence interval nor
a certified total-error bar. All quantities are in dimensionless model units.}
\label{fig:ou-source-work}
\end{figure}

\paragraph{One exterior-domain sensitivity check.}
A separate calculation at fixed $h=0.04$, $\Delta t=0.0025$ and $T=3$
extended $R=20$ to $R=30$, retaining byte-identical source values and
interior coefficients and comparing with the stored $R=20$ calculation.
Across the 31 shared sample times, the absolute-weighted majorant for the
change in the same-method OU-minus-$n=0$ source-work contrast was
$1.14\times10^{-19}$, below the predefined diagnostic threshold
$1.27\times10^{-9}$. This is agreement at floating-point roundoff scale,
not absolute accuracy at $10^{-19}$; common spatial, temporal and profile
errors can remain. The check does not establish an infinite-domain limit,
long-time behavior or control between samples, and neither replaces the
earlier Wiener baseline nor reduces the original comparison guard.

\paragraph{Signed contributions from source-time separations.}
A separate response-kernel calculation uses the fixed $h=0.04$, $R=20$
spatial operator and both genuine quadratures of each source profile.
After averaging the uniform initial phase, its mean source work is
$m[C]=\int_0^1 W(u)C(u)\,du$, where $u$ is the separation of two source
times, not the observation time. The kernel can have either sign, with
$W(0)=2B_{\mathrm{src}}$. On the finest of three lag meshes
($400$, $800$, $1600$ subdivisions), with independent overlap quadratures
of orders $16$ and $32$, the four prescribed quarter-interval contributions
to $\int_0^1 W(u)[C_{\mathrm{OU}}(u)-e^{-u}]\,du$ are, in units of $10^{-6}$,
\[
 (+23.13074873,\;-9.08245885,\;-11.01641674,\;-1.03681136).
\]
Their signed sum is $1.99506177879\times10^{-6}$, with summed heuristic
refinement envelopes $6.98\times10^{-14}$; these are not a certified
error bound or confidence interval. The shortest-lag contribution exceeds
the three negative later-lag contributions (Fig.~\ref{fig:signed-lag-kernel}).
Closure to the stored same-grid Gauss OU and Wiener-rate works passes the
predeclared diagnostic allowances. The separate shift of the Gauss
Wiener-rate work relative to the old RK Wiener baseline is
$-8.47524\times10^{-10}$; it is not assigned to any lag bin.
This calculation is not the $h=0.02$ comparison above and does not exactly
decompose its mixed-integrator result. A lag bin integrates pairs of source
times, not instantaneous positive or negative power. All separations are
at most one model time unit, below the nominal period $2\pi/\rho\simeq3.44$;
the response involves the full generator, not an isolated BIC or a long
ringdown. No stored-energy, formation or stabilization conclusion follows.

\begin{figure}[htbp]
\centering
\includegraphics[width=.98\linewidth]{signed-lag-kernel.pdf}
\caption{Stored finest-grid signed lag kernel $W(u)$ (left) and four fixed
integrals of $W(u)[C_{\mathrm{OU}}(u)-e^{-u}]$ (right), for the frozen $h=0.04$
model. Lines connect the stored lag nodes without smoothing. The caps and
reported net envelope are heuristic numerical sensitivity measures, too
small to resolve at this scale, not confidence intervals or certified
errors; no continuous kernel error band is asserted. The positive net
source-work difference retains all three negative contributions. These
are source-time-pair contributions, not a time history of instantaneous
power or retained energy.}
\label{fig:signed-lag-kernel}
\end{figure}


Scalar radiation pressure has been studied in a related
one-dimensional Q-ball model~\cite{Ciurla2024}; its rates are not imported
into our three-dimensional problem.

Several Q-balls of the same nonlinear field can influence one another,
but a sum of isolated stationary profiles is not automatically a stationary
multi-centre solution. Separation, relative phase and incident frequencies
must be specified. We have not derived a self-consistent environmental
noise spectrum, an interaction with every particle species, or a mechanism
that maintains these excitations indefinitely.

If an additional dynamical mediator is introduced, its spatially nonlocal
static response must be distinguished from temporal memory. Matching a
homogeneous effective potential does not generally match the nonlocal static
interaction. Such extensions are not included in Eq.~\eqref{eq:action}.

An additional open channel also changes the localization condition. For a
real-frequency mode with independent propagating asymptotic channels, every
propagating tail must vanish separately. A cancellation of the original
outgoing amplitude therefore does not establish an embedded bound state of
the extended system. Nor does a zero of a leading-order source overlap
establish this: the coupled background and mode can change at higher order.
These are conditions for a future coupled calculation, not results for an
extension of the model studied here.

\clearpage
\section{Exploratory concentration and constrained transport}\label{app:formation}
A distinct formation hypothesis is that an initially disordered scalar
field amplifies a band of spatial wavelengths and subsequently fragments.
Condensate fragmentation and a scale selected by growing perturbations
are established questions in other Q-ball models~\cite{KasuyaKawasaki2000}.
They do not establish formation in our initial-data family or select the
BIC frequencies. A future comparison should vary narrow-band and broadband
initial spectra at matched total energy and charge, measure growth and
outgoing flux, and distinguish transient concentrations from persistent
localized charged objects. A connected high-density region in a periodic
box should also be checked for closed paths with nonzero winding around
the box. A region with such a path is not an isolated density island at
that threshold. Absence of winding alone would still not establish vacuum
isolation, stationarity or stability; threshold and resolution dependence
must be checked. Total charge conservation forbids production
of net charge from neutral data; opposite charges or a compensating
background would have to be accounted for. No universal integer multiple
of a fundamental wavelength is derived here.

Even persistence alone can be biased by the preparation. For the potential
in Eq.~\eqref{eq:action}, put $S=|\phi|^2$, $\Pi=\partial_t\phi$,
$Q=2\operatorname{Im}\int\phi^*\Pi\,\dd^3x$ and $q=|Q|>0$.
For finite total energy $E=\int(|\Pi|^2+|\nabla\phi|^2+\mathcal U(S))\dd^3x$,
the elementary identity
\[
 E-q=\int\!\left[|\Pi-\ii\operatorname{sgn}(Q)\phi|^2
       +|\nabla\phi|^2-S^2+\tfrac12S^3\right]\dd^3x
\]
implies $\int S^2\dd^3x\ge q-E$ when $E<q$.
Since $\mathcal U(S)\ge S/2$, it also gives
$\operatorname{ess\,sup} S\ge(q-E)/(2E)>0$.
Thus conserved whole-system energy and charge already preclude uniform
dilution to arbitrarily small amplitude for such initial data. This elementary
bound is not a formation result: it specifies neither a
fixed spatial core nor a single stationary Q-ball. Charge transport,
fragmentation and continued oscillation remain possible. A formation test
must therefore resolve localization, profile and phase dynamics rather
than count survival alone. The bound concerns the closed system, not an
unbalanced subvolume or a dissipative boundary model.

\paragraph{A separate finite-time radial relaxation study.}
A subsequent project study of the same potential evolved four prescribed
spherically symmetric Gaussian initial states to $T=1000$, with an
absorbing exterior. Its predeclared diagnostic compared the magnitude of
the charge inside $r<60$ with the stationary-family charge at the
frequency estimated from the central phase, allowing a relative
difference below $10\%$. This tests charge--frequency consistency,
not equality of the evolving field to a stationary solution.
The matched-width state passed at $T=500$ and the broader state at
$T=1000$. The narrower state failed and retained about $2\%$ of its
initial charge inside the measurement region. A larger state remained
strongly pulsating and also failed this diagnostic. Consequently, the
combined formation criterion fixed before the runs was not met.
The successful initial states already had whole-system $E/|Q|<1$;
they do not demonstrate capture of initially unbound radiation.
Two jointly refined spatial and temporal resolutions gave the same
classifications, but absorber reflection was not measured. These reported
finite-time radial observations neither establish exact relaxation nor
test nonspherical stability; a charge-retention plateau alone is not a
binding proof. The original plan and result are identified in the
reproducibility inventory. They have received a document-scope review,
not an independent reproduction of the trajectories.

% Insert after the formation/non-dilution discussion, before Remaining tasks.
% Exploratory numerical diagnostic; distinct from the certified BIC background.
\paragraph{One angular direction during chirped concentration.}
To examine a freedom excluded by radial evolution, we evolved one prescribed
quadrupolar tangent on a time-dependent inward chirp. This is an exploratory
formation diagnostic on a different background from the stationary BIC
studied above. In the units of Eq.~\eqref{eq:action}, the radial initial data are
$F(r)=(3/5)^{3/2}\exp(-r^2/50)\exp(\ii\beta r^2)$,
$\phi_b(r,0)=F(r)$ and $\partial_t\phi_b(r,0)=0.8\ii F(r)$,
with $\beta=0.015$. Their whole-space analytic values are
$Q\simeq240.55177$ and $E/Q\simeq1.038669965>1$; hence the preceding
$E<|Q|$ non-dilution argument does not force a remnant. The calculation
uses a radial Dirichlet box $R=80$ and ends at $T=32$.

For $\delta\phi=\eta(r,t)Y_{20}$, with real unit-normalized $Y_{20}$,
put $w=r\eta$ and $S_b=|\phi_b|^2$. The laboratory-frame tangent equation is
\begin{equation}
 w_{tt}=w_{rr}-\frac{6w}{r^2}
 -(1-4S_b+\tfrac92 S_b^2)w
 -(-2+3S_b)\phi_b^2\overline w .
 \label{eq:angular-chirp-tangent}
\end{equation}
We fix $w(r,0)=r^3(-1/25+2\ii\beta)F(r)$ and
$w_t(r,0)=0.8\ii w(r,0)$, the regular $O(r^3)$ quadrupolar deformation.
The nonlinear radial background is advanced simultaneously, without angular
backreaction. Its stored diagnostics reproduce the earlier chirp evolution;
this replay is a control, not independent formation evidence.

Centered second differences and RK4 are compared at the fixed choices
\[
(h,\Delta t)=(0.1,0.02),(0.05,0.02),(0.1,0.01).
\]
Origin polynomial, independent free-evolution and conjugate-coupling controls
precede the calculation. Our positive reduced diagnostic is
\begin{equation}
 N_h^2=h\sum_{j=1}^{J-1}
 \left[|w_{t,j}|^2+\left(1+\frac6{r_j^2}\right)|w_j|^2\right]
 +\frac1h\sum_{j=0}^{J-1}|w_{j+1}-w_j|^2,
 \qquad r_J=80 .
 \label{eq:angular-chirp-norm}
\end{equation}
It is not a conserved perturbation energy. Positive core, exterior and
collar contributions use fixed intervals $[0,6]$, $[6,76]$ and $[76,80]$,
with shared-node half weights and complete edge contributions.

On the finer spatial grid, the largest \emph{sampled} $N_h(t)/N_h(0)$ is
$1.356829$ at $t=21.6$; its value at $T=32$ is $1.139233$.
Samples are separated by $0.8$, so this is not a continuous-time supremum.
The deformation concentrates in the core and subsequently redistributes
outward (Fig.~\ref{fig:angular-chirp-l2}). Common-grid endpoint differences,
measured in Eq.~\eqref{eq:angular-chirp-norm} and divided by the base-grid
initial norm, are $0.0062849$ for spatial refinement and
$4.0021\times10^{-7}$ for timestep halving. These are sensitivity measures,
not certified error bounds. The sampled collar contribution satisfies
$N_{\mathrm{collar}}^2(t)/N_h^2(0)\le10^{-6}$, the prescribed
contamination threshold. This single direction shows transient amplification,
not a stability classification: the background remains nonstationary,
and other radial shapes, angular sectors, finite amplitudes and later times
remain untested.

\begin{figure}[htbp]
\centering
\includegraphics[width=.98\linewidth]{angular-chirp-l2.pdf}
\caption{Angular response of the inward chirp. Top: the actual signed linear
 density variation $\delta S=2\operatorname{Re}(\overline{\phi_b}\eta)Y_{20}$
 from fine-grid snapshots at $t=0,8,16,24,32$, in the fixed $xz$ window
 $[-12,12]^2$. Radial nodal values are interpolated for rendering; the common
 symmetric colour scale is per unit tangent amplitude. Bottom: global norm
 for all three discretizations, and fine-grid core/exterior squared-norm
 contributions divided by the \emph{global initial} $N_h^2(0)$.
 These are linear tangent data, not a finite-amplitude nonlinear
 three-dimensional evolution or a physical energy partition.}
\label{fig:angular-chirp-l2}
\end{figure}

\paragraph{Leading angular feedback on core-charge transport.}
A separate second-variation calculation follows
$\phi_\varepsilon=B+\varepsilon\eta Y_{20}+\varepsilon^2\Xi+O(\varepsilon^3)$
and evolves the forced spherical mean $\zeta=\langle\Xi\rangle_\Omega$.
Here $\Xi$ is half the second derivative, and $\zeta$ is an angular mean,
not the coefficient of normalized $Y_{00}$. Only this mean is needed
alongside the first tangent for the angle-integrated charge through
second order. The initial family is the volume-preserving ellipsoidal
deformation $r_\varepsilon^2=e^{-2c\varepsilon}(x^2+y^2)+e^{4c\varepsilon}z^2$,
$c=\sqrt{5/(16\pi)}$, with an explicitly differentiated finite-grid
normalization fixing total charge. For $Q_2$, the coefficient of
$\varepsilon^2$, define
\[
 b=\frac{Q_{2,\mathrm{core}}(32)-Q_{2,\mathrm{core}}(0)}{Q_0}
\]
in the same fixed core $r<6$, box $R=80$ and time window.
The three $(h,\Delta t)=(0.1,0.02),(0.05,0.02),(0.1,0.01)$ calculations
give $b=0.08737135,\ 0.08736583,\ 0.08737135$, respectively;
the spatially refined value exceeds the predeclared heuristic guard
$2.864\times10^{-5}$. This guard combines measured spatial, temporal and
charge-balance sensitivities with a fixed floor; it is not a certified
error bound or confidence interval. Stage-integrated physical charge flux
checks the core-charge difference, rather than defining it from the endpoint.

The interpretation requires the initial redistribution
(Fig.~\ref{fig:angular-second-order}). On the finer spatial grid,
$Q_{2,\mathrm{core}}(0)/Q_0=-0.10587389$, whereas
$Q_{2,\mathrm{core}}(32)/Q_0=-0.01850805$.
The radial reference itself has positive net core inflow over this window:
its core fraction increases from $0.58949606$ to $0.65267983$.
The positive second-order net-transport coefficient recovers much of an
initial core-charge deficit but leaves the final correction negative;
it does not show more final charge retained than in the radial reference.
The sampled response is nonmonotonic, with a transient positive core
correction before the negative final value; the decision uses the fixed
endpoint $T=32$, not a selected intermediate maximum.
The preparation fixes charge, not energy:
$E(\varepsilon,0)=E_0+\varepsilon^2 E_2+O(\varepsilon^3)$ with
$E_2/E_0=0.0224455$ on that grid. These are second-order Taylor coefficients at zero deformation,
not a finite-amplitude nonlinear three-dimensional evolution, a binding
comparison at equal energy, or evidence of stationary Q-ball formation.
As a subsequent analytical interpretation, the untruncated continuum
family has $E_{\mathrm{grad}}(\varepsilon)=
E_{\mathrm{grad},0}[2e^{-2c\varepsilon}+e^{4c\varepsilon}]/3$,
while its kinetic and potential integrals remain unchanged. Thus its
positive $E_2=4c^2E_{\mathrm{grad},0}$ is a geometric preparation cost;
this continuum identity neither replaces the stored discrete $E_2$ nor
constitutes a new acceptance criterion or determines the transport sign.

\begin{figure}[htbp]
\centering
\includegraphics[width=.98\linewidth]{angular-second-order-transport.pdf}
\caption{Second-order angular response from stored samples. Left:
$Q_{2,\mathrm{core}}(t)/Q_0$ for all three discretizations, with the initial
fine-space correction marked separately. Right: the initial-subtracted
coefficient $b(t)$ and the negative, independently stage-integrated
outward charge flux divided by $Q_0$, on the fine-space grid.
These routes check the same charge balance, not separate
physical replications. Straight lines join 41 stored samples; no finite
$\varepsilon$ is reconstructed. The final core correction remains negative
despite positive net transport. The endpoint guard is heuristic, not a
confidence interval, certified error or time-uniform uncertainty band.}
\label{fig:angular-second-order}
\end{figure}

% Numerical result independently reviewed: RESULT-REVIEW.txt cee669eda28b331f143750ed01913dc63e985b13ace98badac8e0c3458a652d3.
\paragraph{Compensating the initial energy through second order.}
A further, separately specified family keeps the same first angular
variation but replaces the initial chirp by
$\beta(\varepsilon)=\beta_0+\beta_2\varepsilon^2$.
On each grid, $\beta_2=-E_2/(\partial_\beta E_0)$ is determined from the
preparation energy, without using the later transport. The normalized
family fixes total charge exactly and cancels the initial energy change
through order $\varepsilon^2$; it does not assert exactly equal energies
at finite deformation. Only the homogeneous radial chirp tangent is newly
evolved, and its contribution is combined with the stored angular
second variation. The earlier fixed-charge result remains unchanged.

On the spatially refined grid, $\beta_2=-0.0082897517$ and the normalized
net-transport derivative with respect to $\beta$ is $-11.0018872$.
Their product adds $0.0912029133$ to the previous transport coefficient:
$b_{QE}=0.1785687477$, compared with $b_Q=0.0873658344$.
The base and time-refined values are $0.1785384553$ and $0.1785384478$;
the predeclared heuristic guard for the spatially refined transport
coefficient is $1.525216\times10^{-4}$, not a certified error interval.
The initial core correction is unchanged, while the final core coefficient
$c_{QE}=Q_{2,\mathrm{core},QE}(32)/Q_0$ is now $+0.0726948606$,
compared with $-0.0185080526$ for the previous family.
This final coefficient is reported separately from the
net-transport decision; no additional endpoint-success gate was introduced.
The compensation also reduces the initial inward radial current through
its negative chirp correction. Relative to the radial reference, this
energy-compensated family changes geometry and radial phase together;
its transport response cannot isolate geometry alone.
These coefficients concern the fixed endpoint $T=32$ in the same finite
box, not a finite-amplitude deformed solution, stationary formation or binding.
Figure~\ref{fig:fixed-energy-transport} separates the core coefficients
from their initial-subtracted transport for the two preparations.

\begin{figure}[htbp]
\centering
\includegraphics[width=.98\linewidth]{fixed-energy-transport.pdf}
\caption{Stored spatially refined second-variation coefficients with and
without initial-energy compensation. Left: the core-charge correction;
right: its change from the common initial correction. Straight segments
join 41 stored samples, without fitting or a new evolution. The original
family fixes charge; the compensated family also fixes initial energy
through order $\varepsilon^2$ while changing the radial chirp. Neither
curve represents a finite deformation or a percentage increase. Endpoint
caps on the transport panel are heuristic guards, not confidence intervals,
core-coefficient error bars or time-uniform bounds. The sign change of the
final core coefficient is not a binding or stability result.}
\label{fig:fixed-energy-transport}
\end{figure}

% Stored-state local-clock diagnostic, independently reviewed; finite scope only.
\paragraph{A local two-observable time-shift test.}
The energy-compensated response need not describe a new stationary state:
even an increased endpoint core coefficient could reflect a shifted phase
of the transient concentration. A fixed stored-state check at $T=32$
therefore compares two radial observables,
$C=Q_{\mathrm{core}}/Q_0$ and $H=R_{\mathrm{RMS}}^2/6^2$, where $R_{\mathrm{RMS}}^2$
is the second radial moment weighted by $|\phi|^2$ and divided by its
spatial integral. Write their second-order coefficients as $C_2,H_2$.
Here $C_2(32)=c_{QE}$, not the initial-subtracted $b_{QE}$.
A common local shift $t\mapsto t+\varepsilon^2\tau$ would require
$(C_2,H_2)=\tau(\dot C_0,\dot H_0)$, hence
\[
 \mathcal D=\dot C_0 H_2-\dot H_0 C_2=0.
\]
The moment and its derivative were evaluated directly from the stored
Cauchy data, including the norm-denominator variation; no time-series fit
or additional field evolution was used. On the fine spatial grid,
$H_2=-1.0936132$ and $\mathcal D=0.0021220476$, exceeding the predeclared
heuristic sensitivity guard $6.820019\times10^{-5}$ in inverse model-time
units. Thus this pair is inconsistent with one common local clock shift
under that diagnostic rule. The charge-only shift coefficient is
$\tau_C=C_2/\dot C_0=-9.1439597$, after its separate denominator check:
it multiplies $\varepsilon^2$ and is not a measured finite time shift.
This is only a two-observable endpoint comparison, conditional on the
existing variational outputs; it establishes neither persistence of the
core enhancement nor binding, stationary formation or a full-field distance.
The sensitivity guard is not a certified error bound.

% Independently reviewed complete composite; two production assignments retained.
\paragraph{A finite ellipsoidal response at matched initial charge and energy.}
We additionally evolved finite, axisymmetric ellipsoidal preparations at
$\varepsilon=\pm0.02,\pm0.04$, with initial charge and full discrete energy
matched to each discretization's radial reference within the declared root
tolerance. Unlike the preceding second variation, this tests finite initial
data rather than energy compensation only through order $\varepsilon^2$.
At the fixed endpoint $T=32$ define
\[
 D(e)=\frac{C(+e)+C(-e)-2C(0)}{2e^2},\qquad C=Q_{\mathrm{core}}/Q_0.
\]
This is a symmetric final core response, not initial-subtracted net transport.
On the spatially refined grid,
$D(0.02)=0.0727746617$ and $D(0.04)=0.0730141978$, against the new
same-grid, same-timestep variational reference $c_{\mathrm{ref}}=0.0726948531$.
The prescribed heuristic envelope at $e=0.02$ is $1.2055721\times10^{-4}$;
the discrepancy $7.9808621\times10^{-5}$ plus this allowance meets the
predeclared one-percent agreement criterion. This envelope is not a certified
error bound or a confidence interval. The resolved-remainder criterion is
not met, so no measured Taylor order is claimed.

The separate core-fraction differences at $\varepsilon=+0.02,-0.02$ are
$2.8253326\times10^{-5}$ and $2.9966404\times10^{-5}$, respectively,
both positive in the stored endpoint data; the symmetric envelope is not
a separate certified bound for each sign. These are dimensionless charge-fraction
differences, not energy percentages; a positive symmetric coefficient alone
would not establish both signs. Figure~\ref{fig:fa-response} separates them.
The preparation changes geometry and radial chirp together. This finite-time
comparison establishes neither persistent enhancement, stationary formation,
binding nor general three-dimensional stability. The even-$\ell$, $m=0$
evolution excludes non-axisymmetric dynamics.

The complete twenty-trajectory matrix and separate variational reference
were assembled from two bounded production assignments: fourteen completed
trajectories and the reference were inherited unchanged, and six remaining
trajectories were completed separately. The first assignment retains its
\texttt{INCOMPLETE} status; the full matrix, not a favourable subset,
is used in the comparison.

\begin{figure}[htbp]
\centering
\includegraphics[width=.98\linewidth]{fa-response.pdf}
\caption{Finite symmetric endpoint coefficients for four discretizations
(left) and separate signed core-fraction differences on the spatially refined
grid (right). Caps are heuristic sensitivities, not confidence intervals.
Straight segments join stored time samples; no extrapolation or Taylor-order
fit is used. Initial redistribution is retained rather than subtracted.}
\label{fig:fa-response}
\end{figure}

\begin{figure}[htbp]
\centering
\includegraphics[width=.95\linewidth]{fa-fields-t0.pdf}
\caption{Stored initial amplitude (upper row) and relative phase (lower row)
for the radial and $\varepsilon=\pm0.04$ preparations. The meridional cut
uses $r\le12$ inside the computational domain $R=80$; the circle marks the
measurement radius $r=6$. The common amplitude scale is shared with
Fig.~\ref{fig:fa-fields-final}. Relative phase is
$\arg(\phi_\varepsilon\overline{\phi_0})$ in the common initial phase
convention, masked if either amplitude is below $10^{-4}$ of the shared
maximum. The centre $r<0.1$ is masked; grey is not zero amplitude or phase.}
\label{fig:fa-fields-initial}
\end{figure}

\begin{figure}[htbp]
\centering
\includegraphics[width=.95\linewidth]{fa-fields-t32.pdf}
\caption{The same stored-field reconstruction at $T=32$, with the common
scales and masks of Fig.~\ref{fig:fa-fields-initial}. Complex radial mode
coefficients are interpolated before evaluating the angular basis and phase.
These are axisymmetric amplitude/relative-phase cuts, not a general 3D
simulation, an absolute-phase observable, or a binding diagnostic.
Visually similar colours on the full $[-\pi,\pi]$ scale do not establish
exact phase equality.}
\label{fig:fa-fields-final}
\end{figure}

\paragraph{Phase rotation and the fixed-charge kinetic gap.}
A similar phase colour does not test rigid harmonic motion. We evaluated
necessary conditions on the six stored states with $\varepsilon=0,\pm0.04$
at $t=0,32$, using the complete $R=80$ box. For the reduced mode arrays
$U=r\phi$, $V=\partial_t U$, write
$\|X\|_h^2=4\pi h\sum_{j=1}^{J-1}\sum_\ell|X_{j\ell}|^2$,
$I=\|U\|_h^2$, $K=\|V\|_h^2$ and, for $I>0$,
$\omega_{\mathrm{fit}}=Q/(2I)$.
With $A_h(U)$ the actual discrete nonlinear acceleration, define
\[
 r_v=\frac{\|V-\ii\omega_{\mathrm{fit}}U\|_h^2}{K+\omega_{\mathrm{fit}}^2 I},
 \qquad
 r_a=\frac{\|A_h(U)+\omega_{\mathrm{fit}}^2 U\|_h^2}
                 {\|A_h(U)\|_h^2+\omega_{\mathrm{fit}}^4 I}.
\]
Rigid single-frequency motion requires both numerators to vanish. Initially
$r_v$ vanishes by construction, whereas $r_a\simeq0.0830$; the prescribed
phase velocity does not solve the stationary-profile equation. At $T=32$
the three states give $r_v=0.0282474$--$0.0282631$ and
$r_a=0.0586355$--$0.0588005$. These are squared-residual quotients, not
amplitude percentages or kinetic-energy fractions. This completed endpoint
postprocessing involved no further time integration. It provides neither a core-only classification,
a relaxation law, a certified continuum residual bound nor a general 3D
stability result.

The unnormalized velocity residual has a separate exact algebraic meaning.
In the positive canonical kinetic norm, minimizing velocities at fixed
instantaneous $U$ and total $Q$ gives $V_Q=\ii Q U/(2I)$ and
\[
 E_h(U,V)-E_{Q,h}[U]
 =J_v=\|V-\ii\omega_{\mathrm{fit}}U\|_h^2
 =K-\frac{Q^2}{4I},\qquad
 E_{Q,h}[U]=W_h(U)+\frac{Q^2}{4I},
\]
where $W_h$ contains the radial and angular gradient plus potential energy.
This is a comparison at fixed profile, not a physical projection or measured
relaxation: it generally changes regional charge and need not preserve
other conserved quantities. It does not minimize the spatial profile,
show how energy is radiated, or establish binding. In particular, $r_v$
normalizes $J_v$ by $K+Q^2/(4I)$, not by kinetic or total energy.

\clearpage
% Reviewed stored-data diagnostic; no new PDE evolution or scattering theorem.
% Stored-data diagnostic, not a new PDE evolution or scattering theorem.
\paragraph{Direction of the exterior field.}
For the same chirp, we analysed the stored radial Cauchy data
$u=r\phi_b$, $v=\partial_tu$ at $t=16,24,32$ on all three discretizations.
A fixed window $\chi$ rises from zero to one on $12<r<16$ as
$10s^3-15s^4+6s^5$, $s=(r-12)/4$; $(\chi u,\chi v)$ is zero-padded
onto a periodic analysis interval of length $L=320$. With Fourier series
normalized by $1/M$ for $M$ grid points, both temporal branches are retained:
$a_\sigma=(\widehat{\chi u}+\widehat{\chi v}/(\sigma\ii\Omega))/2$,
$\Omega=\sqrt{1+k^2}$, $\sigma=\pm1$.
The convention $\exp(\ii k r+\sigma\ii\Omega t)$ gives group velocity
$-\sigma k/\Omega$ and branch energy $8\pi L\Omega^2|a_\sigma|^2$.
The denominator $E_{\mathrm{win}}$ below is the \emph{total free spectral window
energy}, including the directionless $k=0$ and unresolved Nyquist bins;
it is neither the original nonlinear energy nor its escaped fraction.

The fine spatial grid gives the values below. Define
$\epsilon_{\mathrm{nl}}=\|\chi(-2S_b+\tfrac32 S_b^2)u\|_2/\|\chi u\|_2$
and $\epsilon_\chi=\|-2\chi'u_r-\chi''u\|_2/\|\chi u\|_2$,
using the radial discrete $L^2$ norm. These are operator-source quotients
in the mass-$1$ units of the model (dimensionally mass squared), not
relative energy errors.
\begin{center}
\begin{tabular}{rrrr}
\toprule
$t$ & $E_{\mathrm{out}}/E_{\mathrm{win}}$ & $\epsilon_{\mathrm{nl}}$ & $\epsilon_\chi$\\
\midrule
16 & 0.904702 & $4.55352\times10^{-5}$ & 0.593850\\
24 & 0.888886 & $1.14544\times10^{-4}$ & 0.828330\\
32 & 0.939306 & $1.58566\times10^{-3}$ & 0.573484\\
\bottomrule
\end{tabular}
\end{center}

All nine decompositions have $E_{\mathrm{out}}>E_{\mathrm{in}}+E_{\mathrm{Nyquist}}$
and $E_{\mathrm{out}}>E_{\mathrm{win}}/2$. The largest change in the outgoing
fraction under spatial refinement is $4.59\times10^{-4}$, and under
timestep halving $1.78\times10^{-8}$; these are sensitivity bounds over
the sampled cases, not certified errors. The exterior nonlinear source
is weak in these units, but the filter source is not: its quotient ranges
from $0.57$ to $0.83$. Zero extension also produces the separately monitored
boundary term $u_r(80^-)\delta(r-80)$, rather than another volume norm.
Thus the fixed-window data are predominantly outward-oriented; without
control of the accumulated source terms, this does not establish
asymptotically free radiation, stationary Q-ball formation or binding of
a remaining core.


\paragraph{A prescribed weak-potential extension.}
In an exploratory, prescribed instantaneous Klein--Gordon--Poisson closure,
matching candidates associated with the first two ladder labels persisted
at two weak couplings each and shifted to lower $\omega^2$. These are four
tested points, not a uniform robustness region. The observation remains
conditional on unresolved Poisson-residual and matching-map regularity
checks; it does not establish robustness in the Einstein--Klein--Gordon
system or nonlinear stability. Including the potential perturbation here
means response within the same instantaneous closure, not full general
relativity. The quantity $2|\Phi(0)|$ denotes twice the central potential
depth and is not identified with the compactness $2GM/R$.


\clearpage
\section{Reproducibility and provenance}\label{app:provenance}
The spatial and phase figures use the stored T2 reconstruction
\path{dipole-forced/reconstruction/mode-1.npz}, the C0 entry of
\path{RUNDE-09/krein1/profile/PROFILE.json}, and the archived
\path{resonance-20260930/winding/RESULT.json}.
The full paths, source hashes, interpolation rules, display thresholds,
rendering code and checks are recorded in
\path{coordination/resonance-20260930/paper-v11-figures/}.
No new field evolution was performed for these figures.
The exploratory two-dimensional comparison is recorded in
\path{coordination/runden-v3/RUNDE-12/leiter2d-praez/}; the additional-channel
probe is in \path{coordination/runden-v3/RUNDE-12/qk1/}.
Both are single-house numerical results with subsequent scope reviews,
not new interval existence certificates. The numerical-lattice comparison
is documented in \path{RUNDE-23/stille-auf-gitter/} and
\path{RUNDE-24/stille-gitter-2/}, relative to the same rounds directory;
the bounded report, symbol and data review is in
\path{coordination/resonance-20260930/grid-robustness-review-20261003/}.
This review does not constitute an independent rerun of the grid spectra.
The prescribed dipole pulse and
its detuning bound are documented in
\path{coordination/resonance-20260930/environment-source-overlap/},
including the result, non-author review and detuning derivation.
The conditional monopole comparison is recorded in
\path{dipole-forced/forced/interior-error/error-budget-components/},
relative to the nonlinear-followup directory below. Its result and
non-author review bind the full input hashes and directed arithmetic.
\path{BEWEISKETTE-KOMPONENTEN.txt} in the parent directory distinguishes
inherited existence and span arguments from the new numerical bounds
and the additional continuation hypotheses. This is a project-internal
review of the assembled argument, not an independent interval
reimplementation of its entire dependency chain.

The $\beta=1$ check is documented under
\path{coordination/resonance-20260930/}, in the dossier
\path{planar-wall-beta1-20261003/}. It includes the fixed local search
plan, four grid/domain stages, stored
complex transmission samples, and non-author code and result reviews.
It is a separate modified-potential calculation, not an additional
existence certificate for the main $\beta=1/2$ model.

The accompanying source manifest records the versions read for this
manuscript. Proof packages retain their own original code hashes,
parameter boxes, interval exports and failed repetitions. The manuscript
assembly neither reruns nor relabels those calculations.
The first proof package covers radial location C0. The second
covers radial location C1 and dipole location C2. Its original
dipole repetition failed a sufficient enclosure test; the subsequently
expanded parameter box is a separately recorded run, not a retroactively
successful preregistered repetition. The primary successful dipole
certificate is logically sufficient on its own.

\begin{longtable}{p{.25\linewidth}p{.67\linewidth}}
\toprule Material&Project-relative source location\\\midrule\endhead
Radial certificate&\path{coordination/runden-v3/RUNDE-07/beweis/}\\
Additional certificates&\path{coordination/runden-v3/RUNDE-10/beweis2/}\\
Cross-house source reviews&\path{coordination/resonance-20260930/beweis1-fremdlesung/}; \path{coordination/resonance-20260930/beweis2-fremdlesung/}\\
Sequence and predictions&\path{coordination/runden-v3/RUNDE-09.md}; \path{coordination/runden-v3/RUNDE-09/VORHERSAGEN-MOD2.md}\\
R13 numerical refinement&\path{coordination/runden-v3/RUNDE-13/leiter3d-praez/ERGEBNIS.md}; \path{coordination/resonance-20260930/paper-r13-review/REVIEW.txt}\\
Spatiotemporal pulse comparison&\path{coordination/resonance-20260930/source-shaping-20261001/}; result and independent result review\\
Factorial actuator comparison&\path{coordination/resonance-20260930/source-shaping-20261001/factorial/}; \path{RESULT.json} and \path{RESULT-REVIEW.txt}\\
Constant pulse-phase comparison&\path{coordination/resonance-20260930/source-shaping-20261001/phase-quadrature/}; result/review and \path{analysis/figure/SUMMARY.json}\\
Stored mode-defect attribution&\path{coordination/resonance-20260930/source-shaping-20261001/mode-defect-audit/}; result, independent review and proposed insertion\\
Quadratic endpoint energy&\path{coordination/resonance-20260930/source-shaping-20261001/phase-quadrature/endpoint-energy/}; \path{RESULT.json} and \path{RESULT-REVIEW.txt}\\
Time-integrated power balance&\path{coordination/resonance-20260930/source-shaping-20261001/phase-quadrature/endpoint-energy/power-balance/}; \path{RESULT.json} and \path{RESULT-REVIEW.txt}\\
Phase-dependent source work&\path{coordination/resonance-20260930/source-shaping-20261001/phase-quadrature/endpoint-energy/power-balance/phase-work/}; coefficient, interval and visualization receipts with reviews\\
Prescribed temporal phase drift&\path{coordination/resonance-20260930/source-shaping-20261001/phase-quadrature/endpoint-energy/power-balance/phase-work/detuning/}; \path{out/RESULT.json}, fixed plan and non-author result review\\
Directional derivative and mirror&The preceding detuning directory, \path{directional/}; \path{out/RESULT.json}, tangent/mirror plan and non-author result review\\
Prescribed phase diffusion&\path{coordination/resonance-20260930/source-shaping-20261001/phase-diffusion/}; \path{out/RESULT.json}, kernel derivations, fixed plan and non-author reviews\\
Rotating mean and finite correlation&The preceding directory, \path{rotating-mean/} and \path{finite-correlation/}; scalar extraction, analytical remainder and failed numerical-acceptance receipt with reviews\\
Gauss finite-correlation comparison&Relative to the phase-diffusion directory above: \path{finite-correlation/gauss-physical/}; \path{out/RESULT.json} and \path{RESULT-REVIEW.txt}, for the separately accepted finite-domain source-work comparison\\
Exterior-domain sensitivity&Relative to the preceding Gauss directory: \path{exterior-sensitivity/}; \path{out/RESULT.json} and \path{RESULT-REVIEW.txt}, for the fixed-time, fixed-resolution domain comparison\\
Signed source-time contributions&Relative to the same Gauss directory: \path{kernel-mechanism/}; \path{out/RESULT.json} and \path{RESULT-REVIEW.txt}, for the signed lag-kernel diagnostic\\
Angular chirp diagnostic&\path{coordination/resonance-20260930/formation-nonlinear-3d/angular-chirp-l2/}; \path{prod/RESULT.json}, non-author result review and visualization receipt\\
Exterior window diagnostic&\path{coordination/resonance-20260930/formation-nonlinear-3d/angular-chirp-l2/exterior-analysis/}; \path{out/RESULT.json}, non-author result review and fixed-window specification\\
Thin-wall derivation&\path{coordination/runden-v3/RUNDE-09/MODELL-DUENNWAND.md}\\
Independent planar-wall scattering&\path{coordination/resonance-20260930/planar-wall-blind-20261002/}; \path{BLIND-ENDE.sha256}, \path{out-v2/RESULT.json}, \path{ERROR-RESULT.json} and non-author review; comparison in \path{coordination/runden-v3/RUNDE-23/tropfen-leiter/ERGEBNIS.md}\\
Finite ellipsoidal response&\path{coordination/resonance-20260930/formation-next-20261002/fixed-energy/finite-amplitude/}; original and rest production records, composite result, independent result review and rendering receipts\\
Phase/profile diagnostics&The preceding directory, \path{rigidity-execution/out/RESULT.json}, \path{PHASE-RIGIDITY-RESULT-REVIEW.txt} and \path{paper-prep/KINETIC-GAP-NOTE.txt}; stored endpoint measurements and separate fixed-charge algebra\\
Radial Gaussian relaxation&\path{coordination/runden-v3/RUNDE-23/bildung-3d/}; \path{KARTE.md}, frozen plan \path{PLAN.md.eingefroren-20261002-220532}, \path{ERGEBNIS.md}; document-scope review in \path{coordination/resonance-20260930/paper-v38-literature/BILDUNG-3D-READING.txt}\\
Evidence audit&\path{coordination/resonance-20260930/novelty-audit/LEITER-STATUS.txt}\\
Quadratic-source check&\path{coordination/resonance-20260930/bic-nonlinear-followup/}\\
Dipole input and response&\path{coordination/resonance-20260930/bic-nonlinear-followup/dipole-forced/}\\\bottomrule
\end{longtable}
These are local project identifiers, not publicly accessible archival
citations. Before submission, the underlying supporting files must be
packaged in a durable public supplement with a license and an archival
identifier. Their current local availability must not be described as a
published independent replication.

\section{Separate numerical-method controls}\label{app:method-controls}
The original finite-correlation field calculation remains a numerical
acceptance failure, and no accepted OU mean is inferred from it. A separate
constant-source control exposed a finite-step defect in the RK4 quadratic
balance: an exact finite-step RK4 identity reproduced the stored endpoint and
accumulated-rate discrepancies, within the prescribed floating-point
comparison tolerances, for 27 case--step combinations sampled at 31 times
each. This identifies the numerical mechanism in those controls, not the
complete cause of the earlier spatial-field failure. A subsequent two-stage
Gauss control evaluated nine constant-system cases (three damping rates and
three complex sources) at three time steps. Its recorded acceptance checks
passed, including independent exact-state and directly integrated-rate
references at the finest step, the prescribed error non-growth checks, and
rejection of an intentionally incorrect rate. The conclusion therefore does
not rest solely on Gauss's algebraic quadratic-balance property. These
controls alone do not validate the original physical-pulse field calculation
or decide the OU-versus-Wiener comparison; the subsequent prescribed-pulse
calculation is reported separately in the main text.

A subsequent spatial adapter used a smooth manufactured solution on the actual
finite operator with $R=20$, $h=0.04$ and 499 interior nodes, through $T=1$.
It covered nine damping rates $\lambda=1,11,\ldots,81$ and a separate zero
control at time steps $0.005$, $0.0025$ and $0.00125$.
Independent full-state references and direct integral references for five
original rates and three modified quadratic rates passed their finest-step
accuracy checks and prescribed error non-growth checks, with a non-author
result review supporting the recorded outcome.
This control alone validates only the manufactured solution, not the prescribed
physical pulse, an OU mean, or a continuum limit.

The finite-step comparison and the Gauss control have separate non-author
result reviews. Their fixed plans, results and reviews are recorded under
\path{coordination/resonance-20260930/source-shaping-20261001/phase-diffusion/finite-correlation/diagnostic-step/},
in \path{recurrence-check/} and \path{gauss-control/}, respectively; the spatial
adapter's separate result and review are in \path{gauss-control/spatial-manufactured/}.
The small-correlation-time derivation and its independent paper review are
recorded in \path{gauss-control/}. These local method checks are not
additional physical findings or interval-certified integration errors.
The subsequent physical-pulse comparison, with its own fixed plan, result,
non-author review and figure receipts, is recorded in
\path{coordination/resonance-20260930/source-shaping-20261001/phase-diffusion/finite-correlation/gauss-physical/}.
Its \path{exterior-sensitivity/} subdirectory records the separate $R=30$
comparison, fixed decision rule, retained $R=20$ reference and non-author
result review. No historical trajectory was rerun for that comparison.
The accepted signed-lag calculation, its fixed plan, independent result
review and stored-data figure are recorded in \path{kernel-mechanism/},
another subdirectory of \path{gauss-physical/}, with
\path{out/RESULT.json}, \path{RESULT-REVIEW.txt} and \path{figure/out/}.
The restricted weak-potential discussion follows the project model review
and its independent algebraic reading in
\path{coordination/resonance-20260930/qstern-paper-review/}, namely
\path{REVIEW-ROOT.txt} and \path{REVIEW-MODEL-NONAUTHOR.txt}; the latter
is not a new audit of the numerical matching claims.
The angular second-variation derivation, fixed plan, code and result reviews
are in \path{coordination/resonance-20260930/formation-next-20261002/};
the accepted result is \path{execution-r2/out/RESULT.json}, and the
stored-data figure is in \path{figure/out/}. The preceding import-only
adapter failure remains recorded separately. The subsequent continuum
preparation-energy identity and its algebraic review are explanatory
checks, not additional prospective gates or physical replications.
The separately specified energy-compensated chirp tangent, its derivation,
plan and code are in the \path{fixed-energy/} subdirectory of
\path{formation-next-20261002/}, with the new stored result in
\path{execution/out/RESULT.json}. Its reference angular second variation
retains its original result and provenance; it was not evolved again.
The energy-compensated comparison figure uses the stored data and receipt
in \path{fixed-energy/figure/out/}. Its separate local-clock check is in
\path{fixed-energy/time-shift/}, including the derivation, fixed endpoint
plan and \path{execution/out/RESULT.json}; it reuses the existing Cauchy
states rather than evolving another trajectory.


\section*{Contributions and disclosure}
\addcontentsline{toc}{section}{Contributions and disclosure}
The working draft was assembled with OpenAI agents from computational and
Anthropic-authored thin-wall section and proof packages. AI-assisted
source reviews are project-internal checks, not journal peer review.
Authorship is assigned to Finn Malte Hinrichsen, Hamburg, at his explicit
direction. No institutional affiliation is asserted. Submission and public
release remain separate decisions.
\begin{thebibliography}{9}

\bibitem{Coleman1985}
S.~Coleman,
``Q-balls,''
\emph{Nuclear Physics B} \textbf{262}, 263--283 (1985),
\href{https://doi.org/10.1016/0550-3213(85)90286-X}{doi:10.1016/0550-3213(85)90286-X};
erratum, \emph{Nuclear Physics B} \textbf{269}, 744 (1986),
\href{https://doi.org/10.1016/0550-3213(86)90520-1}{doi:10.1016/0550-3213(86)90520-1}.

\bibitem{Kovtun2018}
A.~Kovtun, E.~Nugaev, and A.~Shkerin,
``Vibrational modes of Q-balls,''
\emph{Physical Review D} \textbf{98}, 096016 (2018),
\href{https://doi.org/10.1103/PhysRevD.98.096016}{doi:10.1103/PhysRevD.98.096016},
\href{https://arxiv.org/abs/1805.03518v2}{arXiv:1805.03518v2}.

\bibitem{Ciurla2024}
D.~Ciurla, P.~Dorey, T.~Roma\'nczukiewicz, and Y.~Shnir,
``Perturbations of Q-balls: from spectral structure to radiation pressure,''
\emph{Journal of High Energy Physics} \textbf{07} (2024), 196,
\href{https://doi.org/10.1007/JHEP07(2024)196}{doi:10.1007/JHEP07(2024)196},
\href{https://arxiv.org/abs/2405.06591v2}{arXiv:2405.06591v2}.

\bibitem{FriedrichWintgen1985}
H.~Friedrich and D.~Wintgen,
``Interfering resonances and bound states in the continuum,''
\emph{Physical Review A} \textbf{32}, 3231--3242 (1985),
\href{https://doi.org/10.1103/PhysRevA.32.3231}{doi:10.1103/PhysRevA.32.3231}.

\bibitem{YuLu2025}
X.~Yu and Y.~Y.~Lu,
``Existence of Friedrich--Wintgen bound states in the continuum: system of Schr\"odinger equations,''
\href{https://arxiv.org/abs/2504.19573v1}{arXiv:2504.19573v1} (2025).

\bibitem{KasuyaKawasaki2000}
S.~Kasuya and M.~Kawasaki,
``Q-ball formation through the Affleck--Dine mechanism,''
\emph{Physical Review D} \textbf{61}, 041301 (2000),
\href{https://doi.org/10.1103/PhysRevD.61.041301}{doi:10.1103/PhysRevD.61.041301},
\href{https://arxiv.org/abs/hep-ph/9909509v3}{arXiv:hep-ph/9909509v3}.

\bibitem{Azatov2024}
A.~Azatov, Q.~T.~Ho, and M.~M.~Khalil,
``Q-ball perturbations with more details: linear analysis vs lattice,''
\href{https://arxiv.org/abs/2412.13885v1}{arXiv:2412.13885v1} (2024).

\bibitem{ZhangZhouZhu2025}
G.-D.~Zhang, S.-Y.~Zhou, and M.-F.~Zhu,
``$Q$-ball superradiance: Analytical approach,''
\href{https://arxiv.org/abs/2510.27064v1}{arXiv:2510.27064v1} (2025).

\bibitem{ZhangOscillons2020}
H.-Y.~Zhang, M.~A.~Amin, E.~J.~Copeland, P.~M.~Saffin, and K.~D.~Lozanov,
``Classical Decay Rates of Oscillons,''
\href{https://arxiv.org/abs/2004.01202v2}{arXiv:2004.01202v2} (2020),
\href{https://doi.org/10.1088/1475-7516/2020/07/055}{doi:10.1088/1475-7516/2020/07/055}.

\bibitem{Evslin2026}
J.~Evslin, H.~Liu, T.~Roma\'nczukiewicz, Y.~Shnir,
A.~Wereszczy\'nski, and P.~Ziobro,
``Linearized Q-Ball Perturbations,''
\href{https://arxiv.org/abs/2604.07713v1}{arXiv:2604.07713v1} (2026).

\bibitem{Malomed2005}
B.~A.~Malomed, T.~Wagenknecht, A.~R.~Champneys, and M.~J.~Pearce,
``Accumulation of embedded solitons in systems with quadratic nonlinearity,''
\href{https://doi.org/10.1063/1.1938433}{doi:10.1063/1.1938433} (2005),
\href{https://arxiv.org/abs/nlin/0505012}{arXiv:nlin/0505012}.

\bibitem{SofferWeinstein1999}
A.~Soffer and M.~I.~Weinstein,
``Resonances, Radiation Damping and Instability in Hamiltonian Nonlinear Wave Equations,''
\emph{Inventiones mathematicae} \textbf{136}, 9--74 (1999),
\href{https://doi.org/10.1007/s002220050303}{doi:10.1007/s002220050303}.

\bibitem{BattyeSutcliffe2000}
R.~A.~Battye and P.~M.~Sutcliffe,
``Q-ball Dynamics,'' \emph{Nuclear Physics B} \textbf{590}, 329--363 (2000),
\href{https://arxiv.org/abs/hep-th/0003252v1}{arXiv:hep-th/0003252v1}.

\bibitem{Johansson2017Arb}
F.~Johansson, ``Arb: Efficient Arbitrary-Precision Midpoint-Radius
Interval Arithmetic,'' \emph{IEEE Transactions on Computers}
\textbf{66}(8), 1281--1292 (2017),
\href{https://doi.org/10.1109/TC.2017.2690633}{doi:10.1109/TC.2017.2690633}.

\bibitem{PythonFlintDocs}
The Python--FLINT developers, \emph{Python--FLINT 0.9.0 documentation},
\url{https://python-flint.readthedocs.io/en/latest/}
(accessed 1 October 2026).

\bibitem{FlintSoftware}
The FLINT developers, \emph{FLINT: Fast Library for Number Theory},
official repository, \url{https://github.com/flintlib/flint}
(accessed 1 October 2026).

\bibitem{Krawczyk1969}
R.~Krawczyk, ``Newton-Algorithmen zur Bestimmung von Nullstellen mit
Fehlerschranken,'' \emph{Computing} \textbf{4}, 187--201 (1969),
\href{https://doi.org/10.1007/BF02234767}{doi:10.1007/BF02234767}.

\bibitem{Koshelev2018}
K. Koshelev, S. Lepeshov, M. Liu, A. Bogdanov and Y. Kivshar,
``Asymmetric metasurfaces with high-Q resonances governed by bound states
in the continuum,'' \emph{Physical Review Letters} \textbf{121}, 193903 (2018),
\href{https://arxiv.org/abs/1809.00330}{arXiv:1809.00330}.

\bibitem{Flach2003}
S.~Flach, A.~E.~Miroshnichenko, V.~Fleurov, and M.~V.~Fistul,
``Fano resonances with discrete breathers,''
Phys. Rev. Lett. \textbf{90}, 084101 (2003), arXiv:cond-mat/0211313.

\bibitem{Flach2005}
S.~Flach, V.~Fleurov, A.~V.~Gorbach, and A.~E.~Miroshnichenko,
``Resonant light scattering by optical solitons,''
Phys. Rev. Lett. \textbf{95}, 023901 (2005), arXiv:nlin/0409023.

\bibitem{Watabe2012}
S.~Watabe, Y.~Kato, and Y.~Ohashi,
``Excitation Transport through a Domain Wall in a Bose--Einstein Condensate,''
Phys. Rev. A \textbf{86}, 023622 (2012), arXiv:1208.0381.

\bibitem{Heeck2021}
J.~Heeck, A.~Rajaraman, R.~Riley, and C.~B.~Verhaaren,
``Understanding Q-Balls Beyond the Thin-Wall Limit,''
Phys. Rev. D \textbf{103}, 045008 (2021), arXiv:2009.08462.

\bibitem{InagakiMurakami2026}
T.~Inagaki and Y.~Murakami,
``Open-channel radiation zeros and nonlinear damping of a critical-bubble
internal mode,'' arXiv:2609.15056 (2026).

\end{thebibliography}

\end{document}
